%% LyX 2.4.0-alpha3 created this file.  For more info, see https://www.lyx.org/.
%% Do not edit unless you really know what you are doing.
\documentclass[english,tableposition=top]{report}
\usepackage{lmodern}
\renewcommand{\sfdefault}{lmss}
\renewcommand{\ttdefault}{lmtt}
\usepackage[T1]{fontenc}
\usepackage{textcomp}
\usepackage[latin9]{inputenc}
\setcounter{secnumdepth}{3}
\usepackage{color}
\definecolor{shadecolor}{rgb}{0.667969, 1, 1}
\usepackage{babel}
\usepackage{array}
\usepackage{wrapfig}
\usepackage{booktabs}
\usepackage{framed}
\usepackage{url}
\usepackage{amsmath}
\usepackage[unicode=true,pdfusetitle,
 bookmarks=true,bookmarksnumbered=true,bookmarksopen=true,bookmarksopenlevel=2,
 breaklinks=true,pdfborder={0 0 1},backref=section,colorlinks=true]
 {hyperref}

\makeatletter

%%%%%%%%%%%%%%%%%%%%%%%%%%%%%% LyX specific LaTeX commands.
%% Because html converters don't know tabularnewline
\providecommand{\tabularnewline}{\\}

%%%%%%%%%%%%%%%%%%%%%%%%%%%%%% Textclass specific LaTeX commands.
\newenvironment{centred}%
  {\begin{center}\baselineskip=13pt\parskip=1pt}{\end{center}}
\newenvironment{lyxcode}
	{\par\begin{list}{}{
		\setlength{\rightmargin}{\leftmargin}
		\setlength{\listparindent}{0pt}% needed for AMS classes
		\raggedright
		\setlength{\itemsep}{0pt}
		\setlength{\parsep}{0pt}
		\normalfont\ttfamily}%
	 \item[]}
	{\end{list}}

%%%%%%%%%%%%%%%%%%%%%%%%%%%%%% User specified LaTeX commands.
\usepackage{numerica-plus}
\usepackage{upquote}
\newcommand\rel{\,\varrho\;}
\DeclareMathOperator{\erf}{erf}
\DeclareMathOperator{\gd}{gd}

\reuse{}
\constants{ c=30,\omega=0.2 }

\makeatother

\begin{document}
\title{\texttt{numerica-plus}~\\
}
\author{Andrew Parsloe\\
(\url{ajparsloe@gmail.com})}
\maketitle
\begin{abstract}
The \verb`numerica-plus` package defines commands to iterate and
find fixed points of functions of a single variable, to find the zeros
or extrema of such functions, and to calculate the terms of recurrence
relations.\\
\\
\noindent\begin{minipage}[t]{1\columnwidth}%
\begin{shaded}%

\paragraph*{Note:}
\begin{itemize}
\item {\normalsize This document applies to version 2.0.0 of }{\normalsize\texttt{numerica-plus}}{\normalsize .def.}{\small\par}
\item {\normalsize A version of }{\normalsize\texttt{numerica}}{\normalsize{}
from or later than 2021/12/07 is required; (}{\normalsize\texttt{numerica}}{\normalsize{}
requires }{\normalsize\texttt{amsmath}}{\normalsize , }{\normalsize\texttt{mathtools}}{\normalsize{}
and the \LaTeX 3 bundles }{\normalsize\texttt{l3kernel}}{\normalsize{}
and }{\normalsize\texttt{l3packages}}{\normalsize ).}{\small\par}
\item {\normalsize I refer a number of times in this document to }{\normalsize\emph{Handbook
of Mathematical Functions}}{\normalsize , edited by Milton Abramowitz
and Irene A. Segun, Dover, 1965. This is abbreviated to }{\normalsize\emph{HMF}}{\normalsize .}{\small\par}
\item {\normalsize Version 2 of }{\normalsize\texttt{numerica-plus}}{\small\par}
\begin{itemize}
\item {\normalsize is the first stand-alone version; (in v.1 of }{\normalsize\texttt{numerica}}{\normalsize{}
the commands }{\normalsize\texttt{\textbackslash nmcIterate}}{\normalsize ,
}{\normalsize\texttt{\textbackslash nmcSolve}}{\normalsize{} and }{\normalsize\texttt{\textbackslash nmcRecur}}{\normalsize{}
were loaded with the }{\normalsize\texttt{plus}}{\normalsize{} package
option);}{\small\par}
\item {\normalsize makes some small code adjustements;}{\small\par}
\item {\normalsize amends documentation.}{\small\par}
\end{itemize}
\end{itemize}
\end{shaded}%
\end{minipage}
\end{abstract}
\tableofcontents{}

\chapter{Introduction}

Entering 
\begin{lyxcode}
~\textbackslash usepackage\{numerica-plus\}
\end{lyxcode}
in the preamble of your document makes available the commands
\begin{itemize}
\item \verb`\nmcIterate`, a command to iterate a function (apply it repeatedly
to itself), including finding fixed points (values $x$ where $f(x)=x$);
\item \verb`\nmcSolve`, a command to find the zeros of functions of a single
variable (values $x$ for which $f(x)=0$) or, failing that, local
maxima or minima of such functions;
\item \verb`\nmcRecur`, a command to calculate the values of terms in recurrence
relations in a single (recurrence) variable (like the terms of the
Fibonacci sequence or Legendre polynomials).
\end{itemize}
\verb`numerica-plus` requires a version of \verb`numerica` from
or later than 2021/11/26. If found, \verb`numerica` is loaded automatically,
making available the \verb`\nmcEvaluate`, \verb`\nmcInfo`, \verb`\nmcMacros`,
\verb`\nmcConstants`, and \verb`\nmcReuse` commands; see the \verb`numerica`
documentation for details on the use of these commands. 

The commands of the present package all share the syntax of \verb`\nmcEvaluate`.
I will discuss them individually in later chapters but turn first
to something more than a `toy' example that illustrates their use
and gives a sense of `what they are about'.

\section{Example of use: the rotating disk}

\label{sec:introExampleOfUse}Consider a disk rotating uniformly with
angular velocity $\omega$ in an anticlockwise sense in an inertial
system in which the disk's centre \textbf{0} is at rest. Three distinct
points \textbf{1}, \textbf{2}, \textbf{3} are fixed in the disk and,
in a co-rotating polar coordinate system centred at \textbf{0}, have
polar coordinates $(r_{i},\theta_{i})$ ($i,j=1,2,3$). Choose \textbf{01}
as initial line so that $\theta_{1}=0$.

The cosine rule for solving triangles tells us that the time $t_{ij}$
in the underlying inertial system for a signal to pass from \textbf{i}
to \textbf{j} satisfies the equation
\[
t_{ij}=c^{-1}\sqrt{r_{i}^{2}+r_{j}^{2}-2r_{i}r_{j}\cos(\theta_{j}-\theta_{i}+\omega t_{ij})}\equiv f(t_{ij}),
\]
where $c$ is the speed of light. (Equally, we could be describing
an acoustic signal between points on a disk rotating uniformly in
a still, uniform atmosphere \textendash{} in which case $c$ would
be the speed of sound.) Although the equation doesn't solve algebraically
for the time $t_{ij},$ it does tell us that $t=t_{ij}$ is a \emph{fixed
point} of the function $f(t)$. To calculate fixed points we use the
command \verb`\nmcIterate`, or its short-name form \verb`\iter`,
with the star option, \verb`\iter*`. For \verb`\iter` the star option
means: continue iterating until a fixed point has been reached and,
as with the \verb`\eval` command, suppress all elements from the
display save for the numerical result. 

First, though, values need to be assigned to the various parameters.
Suppose we use units in which $c=30,$ and $\omega=0.2$ radians per
second. To avoid having to write these values in the vv-list every
time, I have put in the preamble to this document the statement
\begin{verbatim}
    \constants{ c=30,\omega=0.2 }
\end{verbatim}
For the polar coordinates of \textbf{1 }and \textbf{3 }I have chosen
$r_{1}=10$, $r_{3}=20$ and $\theta_{3}=0.2$ radians (remember $\theta_{1}=0$).
To find a fixed point $t_{13}$ I give $t$ an initial trial value
$1$ (plucked from the air). Its position as the rightmost item in
the vv-list tells \verb`\iter` that $t$ is the iteration variable:
\begin{verbatim}
    \iter*{ c^{-1}\sqrt{r_1^2+r_3^2-2r_1 r_3
              \cos(\theta_3+\omega t)}
          }[ r_1=10,r_3=20,\theta_3=0.2,t=1 ],
    \quad\info{iter}.
\end{verbatim}
$\Longrightarrow$     \iter*{ c^{-1}\sqrt{r_1^2+r_3^2-2r_1 r_3
              \cos(\theta_3+\omega t)}
          }[ r_1=10,r_3=20,\theta_3=0.2,t=1 ],
    \quad\info{iter}. The short-name form of the \verb`\nmcInfo` command from \verb`numerica`
 has been used to display the number of iterations required to attain
the fixed-point value. 

To six figures, only five iterations are needed, which seems rapid
but we can check this by substituting $t=0.356899$ back into the
formula and \verb`\eval`-uating it:
\begin{verbatim}
    \eval*{ c^{-1}\sqrt{r_1^2+r_3^2-2r_1 r_3
              \cos(\theta_3+\omega t)}
          }[ r_1=10,r_3=20,\theta_3=0.2,t=0.356899 ]
\end{verbatim}
$\Longrightarrow$     \eval*{ c^{-1}\sqrt{r_1^2+r_3^2-2r_1 r_3
              \cos(\theta_3+\omega t)}
          }[ r_1=10,r_3=20,\theta_3=0.2,t=0.356899 ], confirming that we have indeed calculated a fixed point. That it
did indeed take only $5$ iterations can be checked by omitting the
asterisk from the \verb`\iter` command and specifying the total number
of iterations to perform. I choose \texttt{do=}7 to show not just
the $5$th iteration but also the next two just to confirm that the
result is stable. We shall view all $7$: \texttt{see=7}. Because
of the length of the formula I have suppressed display of the vv-list
by giving the key \texttt{vvd}\emph{ }an empty value:\footnote{See the associated document \texttt{numerica.pdf}, the chapter on
settings.}
\begin{verbatim}
    \iter[do=7,see=7,vvd=]
      {\[ c^{-1}\sqrt{r_1^2+r_3^2-2r_1 r_3
        \cos(\theta_3+\omega t)} \]}
          [ r_1=10,r_3=20,\theta_3=0.2,t=1 ]
\end{verbatim}
$\Longrightarrow$     \iter[do=7,see=7,vvd=]
      {\[ c^{-1}\sqrt{r_1^2+r_3^2-2r_1 r_3
        \cos(\theta_3+\omega t)} \]}
          [ r_1=10,r_3=20,\theta_3=0.2,t=1 ]
\noindent \begin{flushleft}
The display makes clear that on the $5$th iteration, the $6$-figure
value has been attained. 
\par\end{flushleft}

Alternatively, we could use the \verb`\nmcRecur` command, or its
short-name form \verb`\recur`, to view the successive iterations,
since an iteration is a first-order recurrence: $f_{n+1}=f(f_{n})$:
\begin{verbatim}
    \recur[do=8,see1=0,see2=5,vvd={,\\(vv)\\},*]
      {\[ f_{n+1}=c^{-1}\sqrt{r_1^2+r_3^2-2r_1 r_3
        \cos(\theta_3+\omega f_{n})} \]}
          [ r_1=10,r_3=20,\theta_3=0.2,f_{0}=1 ]
\end{verbatim}
$\Longrightarrow$     \recur[do=8,see1=0,see2=5,vvd={,\\(vv)\\},*]
      {\[ f_{n+1}=c^{-1}\sqrt{r_1^2+r_3^2-2r_1 r_3
        \cos(\theta_3+\omega f_{n})} \]}
          [ r_1=10,r_3=20,\theta_3=0.2,f_{0}=1 ]

\noindent I have specified \texttt{do=8} terms rather than $7$ since
the zero-th term ($f_{0}=1$) is included in the count. I've chosen
to view the last $5$ of them but none prior to those by writing \texttt{see1=0,see2=5}.
Notice the \texttt{vvd} setting, pushing display of the vv-list and
result to new lines and suppressing equation numbering with the \texttt{{*}}
setting (which turns the \verb`multline` environment into a \verb`multline*`
environment).

Another and perhaps more obvious way to find the value of $t_{13}$,
is to look for a zero of the function $f(t)-t$. That means using
the command \verb`\nmcSolve` or its short-name form \verb`\solve`.
I shall do so with the star option \verb`\solve*` which suppresses
display of all but the numerical result. A trial value for $t$ is
required. I have chosen \texttt{t=0}:
\begin{verbatim}
    \solve*{ c^{-1}\sqrt{r_1^2+r_3^2-2r_1 r_3
      \cos(\theta_3+\omega t)} - t }
        [ r_1=10,r_3=20,\theta_3=0.2,t=0 ],
    \quad\nmcInfo{solve}.
\end{verbatim}
$\Longrightarrow$     \solve*{ c^{-1}\sqrt{r_1^2+r_3^2-2r_1 r_3
      \cos(\theta_3+\omega t)} - t }
        [ r_1=10,r_3=20,\theta_3=0.2,t=0 ],
    \quad\nmcInfo{solve}. 

Nearly the same answer as before is attained but this time many more
steps have been required. This is to be expected. The \verb`\solve`
command uses the bisection method. Since $1/2^{10}\approx1/10^{3}$,
about $10$ bisections are needed to determine $3$ decimal places.
Hence we can expect about $20$ bisections for a $6$-decimal-place
answer. The particular form of the \verb`\nmcInfo` command display,
`$1+20$ steps', indicates that it took $1$ search step to find
an interval in which the function vanished and, having found that
interval, $20$ bisections to narrow the position of the zero to $6$-figures.
I will discuss the discrepancy in the final figure in Chapter~\ref{chap:solveSolve};
see �\ref{subsec:solveExtraRounding}.

\subsection{Circuits}

Okay, so we can calculate the time taken in the underlying inertial
system for a signal to pass from one point of the rotating disk to
another. How long does it take to traverse the circuit \textbf{1}
to \textbf{2} to \textbf{3} and back to \textbf{1}? That means forming
the sum $t_{12}+t_{23}+t_{31}$, hence calculating the separate $t_{ij}$
and then using \verb`\eval` to calculate their sum. 

To simplify things, I assume a little symmetry. Let the (polar) coordinates
of \textbf{1} be $(a,0),$ of \textbf{2} be $(r,-\theta)$, and of
\textbf{3} be $(r,\theta)$: \textbf{2} and \textbf{3} are at the
same radial distance from the centre \textbf{0} and at the same angular
distance from the line \textbf{01} but on opposite sides of it, \textbf{3}
ahead of the line, \textbf{2} behind it. The rotation is in the direction
of positive $\theta$. Rather than just calculate $t_{12}+t_{23}+t_{31}$
for the circuit \textbf{1231}, I also calculate the time $t_{13}+t_{32}+t_{21}$
for a signal to traverse the same circuit but in the opposite sense,
\textbf{1321},  and compare them (form the difference). 

Note that with \textbf{2} and \textbf{3} positioned as they are relative
to \textbf{1}, a signal against the rotation from \textbf{3} to \textbf{1}
takes the same time as a signal from \textbf{1} to \textbf{2} and,
in the sense of rotation, a signal from \textbf{2} to \textbf{1} takes
the same time as a signal from \textbf{1} to \textbf{3}. To see this,
suppose the signal from \textbf{2} to \textbf{1} starts at time $t=0$;
it reaches \textbf{1} at a later time $t=t'$ when the disk has rotated
an angle $\omega t'$. Viewed from the underlying inertial system,
the signal path is  a straight line from a point on a circle of radius
$r$ to a point on a concentric circle of radius $a$, the points
subtending an angle at the centre \textbf{0} of $\theta+\omega t'$.
But \textbf{3} at time $t'$ and \textbf{1} at time $t=0$ also subtend
an angle $\theta+\omega t'$ at \textbf{0}, and also lie on circles
of radii $r$ and $a$ respectively. In the underlying inertial system
the line segments \textbf{1}$(0)$\textbf{3}$(t')$ and \textbf{2}$(0)$\textbf{1(}$t')$
are of equal length. Similarly, if a signal from \textbf{3} at time
$t=0$ reaches \textbf{1} at time $t=t''$ then \textbf{3}$(0)$\textbf{1}$(t'')$
and \textbf{1}$(0)$\textbf{2}$(t'')$  are of equal length. Hence
the round trip times are $2t_{12}+t_{23}$ and $2t_{13}+t_{32}$.

\subsubsection{Nesting commands}

Analytically, both $t_{21}$ and $t_{13}$ are the same fixed point
of the function 
\[
c^{-1}\sqrt{r^{2}+a^{2}-2ra\cos(\theta+\omega t)}
\]
 and $t_{31}$ and $t_{12}$ are the same fixed point of the function
\[
c^{-1}\sqrt{r^{2}+a^{2}-2ra\cos(\theta-\omega t).}
\]
To calculate $2t_{12}+t_{23}$ therefore means calculating
\begin{verbatim}
    2\iter*{ c^{-1}\sqrt{a^2+r^2-2ar
                  \cos(\theta-\omega t)} }
      + \iter*{ c^{-1}\sqrt{2r^2-2r^2
                  \cos(2\theta+\omega t)} }
\end{verbatim}
with the analogous expression for $2t_{13}+t_{32}$. But we can do
the comparison of round trip times `in one go' by nesting the \verb`\iter*`
commands inside an \verb`\eval*` command:
\begin{verbatim}
     
    \eval*{ % circuit 1231
        2\iter*[var=t]{ c^{-1}\sqrt{a^2+r^2-2ar
                  \cos(\theta-\omega t)} }[8]
      + \iter*[var=t]{ c^{-1}\sqrt{2r^2-2r^2
                  \cos(2\theta+\omega t)} }[8]
    % circuit 1321
      - 2\iter*[var=t]{ c^{-1}\sqrt{a^2+r^2-2ar
                  \cos(\theta+\omega t)} }[8]
      - \iter*[var=t]{ c^{-1}\sqrt{2r^2-2r^2
                  \cos(2\theta-\omega t)} }[8]
          }[ a=10,r=20,\theta=0.2,t=1 ] 
\end{verbatim}
$\Longrightarrow$     \eval*{ % circuit 1231
        2\times\iter*[var=t]{ c^{-1}\sqrt{a^2+r^2-2ar
                  \cos(\theta-\omega t)} }[8]
      + \iter*[var=t]{ c^{-1}\sqrt{2r^2-2r^2
                  \cos(2\theta+\omega t)} }[8]
    % circuit 1321
      - 2\times\iter*[var=t]{ c^{-1}\sqrt{a^2+r^2-2ar
                  \cos(\theta+\omega t)} }[8]
      - \iter*[var=t]{ c^{-1}\sqrt{2r^2-2r^2
                  \cos(2\theta-\omega t)} }[8]
          }[ a=10,r=20,\theta=0.2,{t}=1 ] . 

By itself this result is of little interest beyond seeing that \verb`numerica-plus`
can handle the calculation. What \emph{is} interesting is to find
values of our parameters for which the time difference vanishes \textendash{}
say values of $\theta$, given the other parameters, especially the
value of $r$. Is there a circuit such that it takes a signal the
same time to travel in opposite senses around the circuit, despite
the rotation of the disk? Rather than nesting the \verb`\iter*` commands
inside an \verb`\eval`, we need to nest them in a \verb`\solve`
command:
\begin{verbatim}
    \solve[p=.,var=\theta,+=1,vvd=\\,*,+=1]
      {\[ % circuit 1231    
        2\times\iter*[var=t,+=1]{ c^{-1}\sqrt{a^2+r^2-2ar
                  \cos(\theta-\omega t)} }
      + \iter*[var=t,+=1]{ c^{-1}\sqrt{2r^2-2r^2
                  \cos(2\theta+\omega t)} }
      % circuit 1321
      - 2\times\iter*[var=t,+=1]{ c^{-1}\sqrt{a^2+r^2-2ar
                  \cos(\theta+\omega t)} }
      - \iter*[var=t,+=1]{ c^{-1}\sqrt{2r^2-2r^2
                  \cos(2\theta-\omega t)} }
      \]}[ a=10,r=20,\theta=0.1,t=1 ]
\end{verbatim}
$\Longrightarrow$     \solve[p=.,var=\theta,+=1,vvd=\\,*]
      {\[ % circuit 1231    
        2\times\iter*[var=t,+=1]{ c^{-1}\sqrt{a^2+r^2-2ar
                  \cos(\theta-\omega t)} }
      + \iter*[var=t,+=1]{ c^{-1}\sqrt{2r^2-2r^2
                  \cos(2\theta+\omega t)} }
      % circuit 1321
      - 2\times\iter*[var=t,+=1]{ c^{-1}\sqrt{a^2+r^2-2ar
                  \cos(\theta+\omega t)} }
      - \iter*[var=t,+=1]{ c^{-1}\sqrt{2r^2-2r^2
                  \cos(2\theta-\omega t)} }
      \]}[ a=10,r=20,\theta=0.1,t=1 ]

One point to note here is the use of \verb`\times` (in \verb`2\times\iter*`).
In this example the formula is displayed (\verb`\solve` wraps around
math delimiters). Without the \verb`\times` the result would have
been the same but the display of the formula would have juxtaposed
the `$2$'s against the following decimals, making it look as if
signal travel times were $20.537778$ and $20.61442$ (and no doubt
causing perplexity). Also note the \texttt{vvd=\textbackslash\textbackslash}
to place the result on a new line and suppress display of the vv-list.

So this expression gives a value of $\theta_{\Delta t=0}$ for one
value of $r$. The obvious next step is to create a table of such
values. I show how that is done in the document \verb`numerica-tables.pdf`
using the command \verb`\nmcTabulate` defined in the associated package
\verb`numerica-tables`. But this is not a research paper on the rotating
disk. I wished to show how the different commands of \verb`numerica-plus`
can be used to explore a meaningful problem. And although it looks
as if a lot of typing is involved, once $c^{-1}\sqrt{r^{2}+a^{2}-2ra\cos(\theta-\omega t)}$
has been formed in \LaTeX{} and values specified in the vv-list, much
of the rest is copy-and-paste with minor editing.

\section{Shared syntax of the new commands}

\texttt{numerica-plus} offers three new commands for three processes:
\verb`\nmcIterate` (short-name form \verb`\iter`) for iterating
functions, \verb`\nmcSolve` (short-name form \verb`\solve`) for
finding the zeros or (local) extrema of functions, and \verb`\nmcRecur`
(short-name form \verb`\recur`) for calculating terms of recurrence
relations.

All three commands share the syntax of the \verb`\nmcEvaluate` (or
\verb`\eval`) command detailed in the associated document \texttt{numerica.pdf}.
When all options are used the command looks like, for instance,
\begin{centred}
\noindent \verb`\nmcIterate*[settings]{expr.}[vv-list][num. format]`
\end{centred}
You can substitute \verb`\nmcSolve`, or \verb`\nmcRecur` for \verb`\nmcIterate`
here. The arguments are similar to those for \verb`\nmcEvaluate`. 
\begin{enumerate}
\item \verb`*` optional switch; if present ensures a single number output
with no formatting, or an appropriate error message if the single
number cannot be produced;
\item \verb`[settings]` optional comma-separated list of \emph{key=value
}settings for this particular command and calculation;
\item \verb`{expr.}` the only mandatory argument; the mathematical expression
in \LaTeX{} form that is the object of interest;
\item \verb`[vv-list]` optional comma-separated list of \emph{variable=value
}items; for \verb`\iter` and \verb`\solve` the \emph{rightmost}
(or innermost) variable in the vv-list may have special significance;
\item \verb`[num. format]` optional format specification for presentation
of the numerical result (rounding, padding with zeros, scientific
notation); boolean output is suppressed for these commands.
\end{enumerate}
Like \verb`\nmcEvaluate`, for all three commands the way the result
is displayed depends on whether the command wraps around math delimiters,
or is used between math delimiters or in the absence of math delimiters.
These distinctions are relevant\emph{ only if the optional star {*}
is absent.}
\begin{itemize}
\item When the star option is used, the \emph{result} is a number only,
without any formatting or vv-list display, or an error message is
displayed.
\item When the star option is not used and one of the following is the case
\begin{itemize}
\item the command wraps around math delimiters, e.g. \verb`\iter{$ expr. $}`,
then 
\begin{itemize}
\item the result is displayed in the form \emph{formula~=~result, (vv-list)}
or the form \emph{formula~$\rightarrow$~result, (vv-list)} as appropriate;
\end{itemize}
\item the command is used within math delimiters, e.g. \verb`\[\iter...\]`,
then 
\begin{itemize}
\item the result is displayed in the form \emph{result, (vv-list)} (without
reference to the formula);
\end{itemize}
\item the command is used in the absence of delimiters, then 
\begin{itemize}
\item the result is presented as if it had been used between \verb`\[`
and \verb`\]`.
\end{itemize}
\end{itemize}
\end{itemize}
Looking at the various examples in the preceding section on the rotating
disk you will see illustrations of all these situations.

\subsection{Settings}

Nearly all the settings available to the \verb`\eval` command are
available to these other commands. To save switching between documents
I reproduce in Table~\ref{tab:introSettingsInherited} the options
found in \texttt{numerica.pdf}, although for discussion of the options
you will need to refer to that document. In addition, each of the
present commands also has settings of its own, discussed at the relevant
parts of the following chapters.

\subsection{Nesting}

In v.1 of \verb`numerica`, for commands to be nested one within an
another, it was necessary for the inner command to be starred (and
thus produce a purely numerical result). With v.2 of \verb`numerica`
this is no longer the case. A nested command is detected as such and
the star automatically set, whether the user has explicitly starred
the command or not. Provided the starred form of a command actually
does produce a numerical result and not an error message then it can
be nested within the main argument of any one of the other commands,
including itself. The example of use, �\ref{sec:introExampleOfUse}
above, shows several examples of this. The starred form can also be
used in the vv-list of any one of the commands, including itself.
The associated document \texttt{numerica.pdf} shows examples of an
\verb`\eval*` command being used in the vv-list of an \verb`\eval`
command. 

\begin{table}
\noindent \centering{}\caption{\protect\label{tab:introSettingsInherited}Inherited settings options}
\noindent \begin{center}
\begin{tabular}{ll>{\raggedright}p{4cm}>{\raggedright}p{4cm}}
\toprule 
{\small key} & {\small type} & {\small meaning} & {\small default}\tabularnewline
\midrule
{\small\texttt{dbg}} & {\small int} & {\small debug `magic' integer} & {\small\texttt{0}}\tabularnewline
{\small\texttt{\textasciicircum}} & {\small char} & {\small exponent mark for sci. notation input} & {\small\texttt{e}}\tabularnewline
{\small\texttt{xx}} & {\small int (0/1)} & {\small multi-token variable switch} & {\small\texttt{1}}\tabularnewline
{\small\texttt{()}} & {\small int (0/1/2)} & {\small trig. function arg. parsing} & {\small\texttt{0}}\tabularnewline
{\small\texttt{o}} &  & {\small degree switch for trig. funcions} & \tabularnewline
{\small\texttt{log}} & {\small num} & {\small base of logarithms for }{\small{\small\verb`\log`}} & {\small\texttt{10}}\tabularnewline
{\small\texttt{vvmode}} & {\small int (0/1)} & {\small vv-list calculation mode} & {\small\texttt{0}}\tabularnewline
{\small\texttt{vvd}} & {\small tokens} & {\small vv-list display-style spec.} & {\small\texttt{\{,\}\textbackslash mskip 12mu plus 6mu minus 9mu(vv)}}\tabularnewline
{\small\texttt{vvi}} & {\small token(s)} & {\small vv-list text-style spec.} & {\small\texttt{\{,\}\textbackslash mskip 36mu minus 24mu(vv)}}\tabularnewline
{*} &  & {\small suppress equation numbering if }{\small\texttt{\textbackslash\textbackslash}}{\small{}
in }{\small\texttt{vvd}} & \tabularnewline
{\small\texttt{p}} & char(s) & {\small punctuation (esp. in display-style)} & {\small\texttt{, (comma)}}\tabularnewline
{\small\texttt{S+}} & {\small int} & {\small extra rounding for stopping criterion for sums} & {\small\texttt{2}}\tabularnewline
{\small\texttt{S?}} & {\small$\text{int}\ge0$} & {\small stopping criterion query terms for sums} & {\small\texttt{0}}\tabularnewline
{\small\texttt{P+}} & {\small int} & {\small extra rounding for stopping criterion for products} & {\small\texttt{2}}\tabularnewline
{\small\texttt{P?}} & {\small$\text{int}\ge0$} & {\small stopping criterion query terms for products} & {\small\texttt{0}}\tabularnewline
\bottomrule
\end{tabular}
\par\end{center}
\end{table}


\chapter{Iterating functions: \texttt{\textbackslash nmcIterate}}

\label{chap:Iterating-functions}Only in desperation would one try
to evaluate a continued fraction by stacking fraction upon fraction
upon fraction like so:
\begin{verbatim}
    \eval{\[ 1+\frac{1}{1+\frac{1}{1+\frac{1}
      {1+\frac{1}{1+\frac{1}{1+\frac{1}
        {1+\frac{1}{1+\frac{1}{1+\frac{1}
          {1+\frac{1}{1+\frac{1}
            {1+\frac{1}{1}}}}}}}}}}}} \]}
\end{verbatim}
$\Longrightarrow$     \eval{\[ 1+\frac{1}{1+\frac{1}{1+\frac{1}
      {1+\frac{1}{1+\frac{1}{1+\frac{1}
        {1+\frac{1}{1+\frac{1}{1+\frac{1}
          {1+\frac{1}{1+\frac{1}
            {1+\frac{1}{1}}}}}}}}}}}} \]}

\noindent \verb`numerica-plus` provides a command for tackling problems
like this sensibly. In such problems a function is repeatedly applied
to itself (iterated). This is done through the command \verb`\nmcIterate`
or (short-name form) \verb`\iter`. Thus to evaluate this continued
fraction write (for instance),
\begin{centred}
\verb`\iter[do=15,see=5]{\[ 1+1/x \]}[x=1]` $\Longrightarrow$ \iter[do=15,see=5]{\[ 1+1/x \]}[x=1]
\end{centred}
The \verb`\iter` command evaluates $1+1/x$ when $x=1$ and then
uses this value as a new $x$-value to substitue into $1+1/x$, to
once again evaluate and use as a new $x$-value, and so on. It looks
as if the repeated iterations are approaching \verb`\eval{$ \tfrac{\sqrt{5}+1}2 $}`
$\Longrightarrow$ \eval{$ \tfrac{\sqrt{5}+1}2 $}. Increasing the
number of iterations in the example from \texttt{do=15} to, say, \texttt{do=18},
shows that this is indeed the case. 
\begin{centred}
\verb`\iter[do=18,see=5]{\[ 1+1/x \]}[x=1]` $\Longrightarrow$ \iter[do=18,see=5]{\[ 1+1/x \]}[x=1] 
\end{centred}
But iteration of functions is not limited to continued fractions.
Particularly since the emergence of chaos theory, iteration has become
an important study in its own right. Any function with range within
its domain can be iterated \textendash{} repeatedly applied to itself
\textendash{} like the cosine:
\begin{centred}
\verb`\iter[do=20,see=4]{\[ \cos x \]}[x=\pi/2]` $\Longrightarrow$
\iter[do=20,see=4]{\[ \cos x \]}[x=\pi/2]
\end{centred}
which displays the first one and last four of 20 iterations of $\cos x$
when $x=\tfrac{\pi}{2}$. It looks as if the cosine is `cautiously'
approaching a limit, perhaps around $0.738$ or $0.739$. You need
to nearly double the number of iterations (\texttt{do=40}) to confirm
that this is so.

The logistic function $kx(1-x)$ exhibits a variety of behaviours
depending on the value of $k$. For instance, with $k=3.5$ we get
a period-4 cycle:
\begin{centred}
\verb`\iter[do=12,see=8]{\[ kx(1-x) \]}[k=3.5,x=0.5]` $\Longrightarrow$
 \iter[do=12,see=8]{\[ kx(1-x) \]}[k=3.5,x=0.5]
\end{centred}
and with $k=3.1$ we get a period-2 cycle, although it takes many
more iterations to stabilize there:
\begin{centred}
\verb`\iter[do=42,see=4]{\[ kx(1-x) \]}[k=3.1,x=0.5]` $\Longrightarrow$
\iter[do=42,see=4]{\[ kx(1-x) \]}[k=3.1,x=0.5]
\end{centred}

\section{Star (\texttt{{*}}) option: fixed points}

In the first two of these examples, iteration eventually ended at
a \emph{fixed point}. This is a point $x$ where $f(x)=x$. Appending
a star (asterisk) to the \verb`\iter` command is the signal for iteration
to continue until a fixed point has been reached at the specified
rounding value:
\begin{centred}
\verb`\iter*{ 1+a/x }[a=n(n+1),n=1,x=1]` $\Longrightarrow$ \iter*{ 1+a/x }[a=n(n+1),n=1,x=1]
\end{centred}
(with the default rounding value $6$).\footnote{For your own interest try also putting $n=2,3,4,\ldots$ in the vv-list
of this expression.} The star overrides any value for the number of iterations to perform
(the \texttt{do} key) that may have been entered in the settings option.
It also overrides any elements of the display other than the numerical
result. With the star option math delimiters are irrelevant \textendash{}
other than displaying minus signs correctly when \verb`\iter*` is
between them.

A function may not approach a fixed point when iterated \textendash{}
see the examples with the logistic function above. To prevent an infinite
loop \verb`\iter*` counts the number of iterations performed and
when that number reaches a certain cut-off value \textendash{} the
default is $100$ \textendash{} the loop terminates and a message
is displayed:
\begin{centred}
\verb`\iter*{kx(1-x)}[k=3.5,x=0.5]` $\Longrightarrow$ \iter*{kx(1-x)}[k=3.5,x=0.5]
\end{centred}
In this case we \emph{know} that a fixed point does not exist, but
that may not always be the case. One response to a message like this
is to change parameter values or starting value of the iteration variable.
For instance, changing the parameter value to $k=1.5$,
\begin{centred}
\verb`\iter*{kx(1-x)}[k=1.5,x=0.5]` $\Longrightarrow$ \iter*{kx(1-x)}[k=1.5,x=0.5],
\end{centred}
means a fixed point is now attained. It is easy to check that $1/3$
is indeed a fixed point (but that makes the $4$ in the last decimal
place a concern; see the extra rounding setting, �\ref{subsec:iterExtra-rounding}). 

But should a fixed point still not eventuate after `fiddling' with
parameter and start values, there are two general adjustments one
might try: either 
\begin{enumerate}
\item reduce the rounding value, from the default $6$ (or the one specified),
to a smaller value, or
\item increase the cut-off figure from the default 100 to some higher value. 
\end{enumerate}
The former is done via the trailing number format optional argument
of the \verb`\iter` command; the latter is done via the settings
option, see �\ref{sec:iterSettings-option}, specifically �\ref{subsec:iterMaximum-iteration-count}. 

\subsection{Use with \texttt{\textbackslash nmcInfo}}

It is of interest to know how many iterations are required to reach
a fixed point at a particular rounding value. That knowledge allows
a good guess as to whether a fixed point will be attained at a greater
rounding value. Thus when iterating the function
\[
f(t_{ij})=c^{-1}\sqrt{r_{i}^{2}+r_{j}^{2}-2r_{i}r_{j}\cos(\theta_{j}-\theta_{i}+\omega t_{ij})}
\]
in �\ref{sec:introExampleOfUse} only $5$ iterations were required
to attain $6$-figure accuracy for the fixed point. That information
came by following the \verb`\iter*` command with \verb`\nmcInfo`
(or \verb`\info`) with the argument \verb`iter`. And generally,
for any `infinite' process, follow the command with an \verb`\info`
command if you want to know how many `steps' \textendash{} in the
present case iterations \textendash{} are required to achieve the
result. So, if $5$ iterations achieve $6$-figure accuracy, presumably
something like $10$ iterations will achieve $12$-figure accuracy:
\begin{verbatim}
    \iter*{ c^{-1}\sqrt{r_i^2+r_j^2-2r_i r_j
      \cos(\theta_{ij}+\omega t)}
        }[ r_i=10,r_j=20,\theta_{ij}=0.2,t=1 ][12]
    ,\quad\info{iter}. 
\end{verbatim}
$\Longrightarrow$     \iter*{ c^{-1}\sqrt{r_i^2+r_j^2-2r_i r_j
      \cos(\theta_{ij}+\omega t)}
        }[ r_i=10,r_j=20,\theta_{ij}=0.2,t=1 ][12]         
    ,\quad\info{iter}. (Remember, \verb`numerica-plus` knows the values of $c$ and $\omega$
from a \verb`\constants` statement in the preamble.) And indeed only
$9$ iterations suffice to achieve $12$-figure accuracy: 
\begin{verbatim}
    \iter[do =11,see=4]
      { c^{-1}\sqrt{r_i^2+r_j^2-2r_i r_j
        \cos(\theta_{ij}+\omega t)}
      }[ r_i=10,r_j=20,\theta_{ij}=0.2,t=1 ][12]
\end{verbatim}
$\Longrightarrow$     \iter[do =11,see=4]
      { c^{-1}\sqrt{r_i^2+r_j^2-2r_i r_j
        \cos(\theta_{ij}+\omega t)}
      }[ r_i=10,r_j=20,\theta_{ij}=0.2,t=1 ][12]

Or again, with another example from earlier,
\begin{centred}
\verb`$ \iter*{\cos x}[x=\pi/2] $,\ \info{iter}.` $\Longrightarrow$
$\iter*{\cos x}[x=\pi/2] $,\ \info{iter}.
\end{centred}
That suggests that around $2\times37=74$ iterations will give a $2\times6=12$-figure
answer, well within the cut-off figure of $100$:
\begin{centred}
\verb`$ \iter*{\cos x}[x=\pi/2][12] $,\ \info{iter}.` $\Longrightarrow$
$\iter*{\cos x}[x=\pi/2][12] $,\ \info{iter}.
\end{centred}

\section{Settings option}

\label{sec:iterSettings-option}The settings option is a comma-separated
list of items of the form \emph{key~=~value}.

\subsection{Inherited settings}

Nearly all of the keys discussed in the settings\emph{ }option for
\verb`\nmcEvaluate` are available for \verb`\nmcIterate`. Table~\ref{tab:introSettingsInherited}
above lists these, repeating a table from \texttt{numerica.pdf}. Thus
should a quantity in the vv-list depend on the iteration variable,
forcing an implicit mode calculation, simply enter, as with \verb`\eval`,
\verb`vv@=1` (alternatively, \verb`vvmode=1`) in the settings option:
\begin{centred}
\verb`\iter*[vv@=1]{$ f(x) $}[f(x)=1+a/x,a=12,x=1]` $\Longrightarrow$
\iter*[vv@=1]{$ f(x) $}[f(x)=1+a/x,a=12,x=1].
\end{centred}
Implicit in the example is the default multi-token setting \texttt{xx=1}
inherited from \verb`\eval` and ensuring that the multi-token variable
$f(x)$ is treated correctly. 

Let's add \verb`dbg=1` to the example:
\begin{centred}
\verb`\iter*[dbg=1,vv@=1]{$ f(x) $}[f(x)=1+a/x,a=12,x=1]` $\Longrightarrow$
\iter*[dbg=1,vv@=1]{$ f(x) $} [f(x)=1+a/x,a=12,x=1]
\end{centred}
The multi-token variable \verb`f(x)` has been changed to a single-token.
The values shown under `stored' and `fp-form' are those of the
\emph{final} iteration.

\subsection{\texttt{\textbackslash nmcIter}ate-specific settings}

In addition to the inherited settings there are some specific to \verb`\nmcIterate`.
These are listed in Table~\ref{tab:iterSettings}.

\subsubsection{Iteration variable}

In nearly all of the examples so far, the iteration variable has been
the rightmost variable in the vv-list and has not needed to be otherwise
specified. However it is sometimes not feasible to indicate the variable
in this way. In that case, entering
\begin{lyxcode}
var~=~<variable~name>
\end{lyxcode}
in the settings option enables the variable to be specified, irrespective
of what the rightmost variable in the vv-list is. Here, \texttt{<variable
name>} will generally be a character like \texttt{x} or \texttt{t}
or a token like \verb`\alpha`, but it could also be a multi-token
name like {\ttfamily\verb`x'`}\texttt{ }or \texttt{\textbackslash beta\_\{ij\}}
(or even \texttt{Fred} if you so chose). Although the iteration variable
can be independently specified like this, it must still be given an
initial \emph{value} in the vv-list \textendash{} only it need not
be the rightmost variable.

In the following example the rightmost variable is $a$ which is clearly
\emph{not} the iteration variable:
\begin{centred}
\verb`\iter[var=x,do=40,see=5]{$ 1+a/x $}[x=a/6,a=6][*]` $\Longrightarrow$
\iter[var=x,do=40,see=5]{$ 1+a/x $}[x=a/6,a=6][*]
\end{centred}

\subsubsection{Extra rounding for fixed-point calculations}

\label{subsec:iterExtra-rounding}
\begin{table}
\centering{}\caption{\protect\label{tab:iterSettings}Settings for \texttt{\textbackslash nmcIterate}}
\begin{center}
\begin{tabular}{llll}
\toprule 
{\small key} & {\small type} & {\small meaning} & {\small default}\tabularnewline
\midrule
{\small\texttt{var}} & {\small token(s)} & {\small iteration variable} & \tabularnewline
{\small\texttt{+}} & {\small int} & {\small fixed point extra rounding} & {\small\texttt{0}}\tabularnewline
{\small\texttt{max}} & {\small int > 0} & {\small max. iteration count (fixed points)} & {\small\texttt{100}}\tabularnewline
{\small\texttt{do}} & {\small int > 0} & {\small number of iterations to perform} & {\small\texttt{5}}\tabularnewline
{\small\texttt{see}} & {\small int > 0} & {\small number of final iterations to view} & {\small\texttt{4}}\tabularnewline
 & {\small int ($\mathtt{0}/\mathtt{1}/\mathtt{2}$)} & {\small form of result saved with }{\small{\small\verb`\`}} & {\small\texttt{0}}\tabularnewline
\bottomrule
\end{tabular}
\par\end{center}
\end{table}
\verb`numerica-plus` determines that a fixed point has been reached
when the difference between successive iterations vanishes when rounded
to the current rounding value. One might want reassurance that this
really is the correct value by seeking a fixed point at a higher rounding
value than that displayed. This extra rounding is achieved by entering 
\begin{lyxcode}
+~=~<integer>
\end{lyxcode}
in the settings option. By default this extra rounding is set to zero. 

We have seen before that $\cos x$ starting at $x=\tfrac{1}{2}\pi$
takes $37$ iterations to reach a $6$-figure fixed point $0.739085$,
about $6$ iterations per decimal place. By entering \texttt{+=1}
in the settings option the number of iterations is increased to $43$,
$6$ more than $37$ but, reassuringly, the $6$-figure result that
is displayed remains unchanged:
\begin{centred}
\verb`$ \iter*[+=1]{\cos x}[x=\pi/2] $,\ \info{iter}.` $\Longrightarrow$
$ \iter*[+=1]{\cos x}[x=\pi/2] $,\ \info{iter}.
\end{centred}

\subsubsection{Maximum {\small iteration count for fixed-point searches}}

\label{subsec:iterMaximum-iteration-count}To prevent a fixed-point
search from continuing indefinitely when no fixed point exists, there
needs to be a maximum number of iterations specified after which point
the search is called off. By default this number is $100$. To change
it enter
\begin{lyxcode}
max~=~<positive~integer>
\end{lyxcode}
in the settings option.

\subsubsection{Number of iterations to perform}

To specify the number of iterations to perform enter
\begin{lyxcode}
do~=~<positive~integer>
\end{lyxcode}
in the settings option. Note that if the \verb`*` option is present
this value will be ignored and iteration will continue until either
a fixed point or the maximum iteration count is reached. By default
\texttt{do} is set to $5$. (Note that \texttt{do} can be set to a
greater number than \texttt{max}; \texttt{max} applies only to \verb`\iter*`.)

\subsubsection{Number of iterations to show}

To specify the number of final iterations to show enter 
\begin{lyxcode}
see~=~<positive~integer>
\end{lyxcode}
in the settings option. By default \texttt{see} is set to $4$. Always
it is the \emph{last} \texttt{see} iterations that are displayed.
If \texttt{see} is set to a greater value than \texttt{do}, all iterations
are shown. If the star option is used the \texttt{see} value is ignored.

\subsubsection{Form of result saved by \texttt{\textbackslash nmcReuse}}

By entering
\begin{lyxcode}
reuse~=~<integer>
\end{lyxcode}
in the settings option of the \verb`iter` command it is possible
to specify the form of result that is saved when using \verb`\nmcReuse`.
(This setting has no effect when the star option is used with \verb`\nmcIterate`.
In that case only the numerical result of the fixed point calculation
\textendash{} if successful \textendash{} is saved.) The possibilities
are:
\begin{itemize}
\item \texttt{int=0} (or any integer $\mathtt{\ne1,2}$) saves the display
resulting from the \verb`\iter` command (the default);
\item \texttt{int=1} saves a comma-separated list of braced pairs of the
form:\texttt{ }~\\
\texttt{\{$k$, value-of-$k$-th-iterate\}};
\item \texttt{int=2} saves a comma-separated list of iterate values.
\end{itemize}
Note that the number and content of the items in the lists are those
resulting from the \texttt{see} setting (the number of iterations
to view).
\begin{verbatim}
    \iter[reuse=1,do=12,see=4]
      {\[ kx(1-x) \]}[k=3.5,x=0.5]
    \reuse{logistic}
\end{verbatim}
$\Longrightarrow$ \iter[reuse=1,do=12,see=4]{\[ kx(1-x) \]}[k=3.5,x=0.5]  \reuse[renew]{logistic}

\noindent whence \verb`\logistic` $\Longrightarrow$ \logistic{}.
As you can see the control sequence \verb`\logistic` displays as
a comma-separated list of numbers, alternating between the iterate
ordinal and the iterate value. That these are stored as braced pairs
can be seen by using \TeX 's \verb`\meaning` command:
\begin{verbatim}
    \meaning \logistic
\end{verbatim}
$\Longrightarrow$ \meaning \logistic

\subsection{Changing default values}

\begin{wraptable}[11]{o}{0.5\columnwidth}%
\centering{}\vspace{-2.7ex}
\caption{\protect\label{tab:iterSettingsDefaults}Defaults for \texttt{\textbackslash nmcIterate}}
\begin{center}
\begin{tabular}{ll}
\toprule 
{\small key} & {\small default}\tabularnewline
\midrule
{\small iter-extra-rounding} & {\small\texttt{0}}\tabularnewline
{\small iter-max-iterations} & {\small\texttt{100}}\tabularnewline
{\small iter-do} & {\small\texttt{5}}\tabularnewline
{\small iter-see-last} & {\small\texttt{4}}\tabularnewline
{\small iter-reuse} & {\small\texttt{0}}\tabularnewline
\bottomrule
\end{tabular}
\par\end{center}\end{wraptable}%
If you wish to change the default values of the various settings for
\verb`\nmcIterate` this can be done by entering new values in a configuration
file \texttt{numerica-plus.cfg} as described in the chapter on settings
in the associated document \texttt{numerica.pdf}. The relevant keys
are listed in Table~\ref{tab:iterSettingsDefaults}, corresponding
to the \verb`+`, \verb`max`, \verb`do`, \verb`see` and \verb`reuse`
settings of the \verb`\iter` command. (Obviously it makes no sense
to have a default setting for the iteration variable. That will change
from case to case.)

\section{Errors}

By errors I refer to \verb`numerica-plus` errors rather than \LaTeX{}
errors. We have already met one in the discussion of fixed points:
\begin{centred}
\verb`\iter*{kx(1-x)}[k=3.5,x=0.5]` $\Longrightarrow$ \iter*{kx(1-x)}[k=3.5,x=0.5]
\end{centred}
For a function to be iterated indefinitely, its range must lie within
or be equal to its domain. If even part of the range of a function
lies outside its domain, then on repeated iteration there is a chance
that a value will eventually be calculated which lies in this `outside'
region. Iteration cannot continue beyond this point and an error message
is generated. As an example consider the inverse cosine, \verb`\arccos`.
This can be iterated only so far as the iterated values lie between
$\pm1$ inclusive. If we try to iterate \verb`\arccos` at 0 for example,
since $\cos\frac{1}{2}\pi=0$, $\arccos0=\eval{0.5\pi}[4]$ (which
is $\tfrac{1}{2}\pi$) so only a first iterate is possible. But we
could choose an initial value more carefully; $37$ iterations of
the cosine at $\tfrac{1}{2}\pi$ led to a fixed point $0.739085$,
so let's choose $0.739085$ as initial point and perform $37$ iterations:
\begin{centred}
\verb`\iter[do=37,see=4]{\[ \arccos x \]}[x=0.739085]` $\Longrightarrow$
\iter[do=37,see=4]{\[ \arccos x \]}[x=0.739085]
\end{centred}
The result of the $37$th iteration is greater than $1$. Thus increasing
the number of iterations to 38 should generate an error message:
\begin{centred}
\verb`\iter[do=38,see=4]{\[ \arccos x \]}[x=0.739085]` $\Longrightarrow$\iter[do=38,see=4]{\[ \arccos x \]} [x=0.739085]  
\end{centred}
\verb`l3fp` objects when asked to find the inverse cosine of a number
greater than $1$.

\chapter{Finding zeros and extrema: \texttt{\textbackslash nmcSolve}}

\label{chap:solveSolve}\verb`numerica-plus` provides a command,\textbf{
}\verb`\nmcSolve` (short-name form \verb`\solve`), for finding a
zero of a function, should it have one. In the following example, 
\begin{centred}
\verb`\solve[p]{\[ e^{ax}-bx^2 \]}[a=2,b=3,{x}=0]` $\Longrightarrow$
\solve[p]{\[ e^{ax}-bx^2 \]}[a=2,b=3,{x}=0]
\end{centred}
I have sought and found a solution $x$ to the equation $e^{ax/2}-bx^{2}=0$
when $a=2$ and $b=3$, starting with a trial value $x=0$, entered
as the \emph{rightmost} variable in the vv-list (and em-braced since
I don't want this trial value displaying in the presentation of the
result). Although $x$ has been found to the default six-figure accuracy,
it is evident that the function vanishes only to five figures. Let's
check:
\begin{centred}
\verb`\eval{$ bx^2 $}[b=3,x=x=-0.390647]` $\Longrightarrow$ \eval{$ bx^2 $}[b=3,x=-0.390647], 

\verb`\eval{$ e^{ax} $}[a=2,x=-0.390647]` $\Longrightarrow$ \eval{$ e^{ax} $}[a=2,x=-0.390647]; 
\end{centred}
the values agree save in the final digit.

This discrepancy in the final decimal place or places is a general
feature of solutions found by \verb`\solve`. It is the value of $x$,
not the value of $f(x)$, that is being found (in this case) to six
figures. If the graph of a function crosses the $x$-axis steeply
then the $x$ value (the zero) may be located to a higher precision
than the function value. Conversely, if the graph of a function crosses
the $x$-axis gently (at a shallow angle) then the function value
will vanish to a greater number of decimal places than the zero (the
$x$ value) is found to.

A second example, which we can check against values tabulated in \emph{HMF},
is to find a value of $x$ that satisfies $\tan x=\lambda x$. In
other words, find a zero of $\tan x-\lambda x$. In the example $\lambda$
is negative, so a trial value for $x$ greater than $\pi/2$ seems
like a good idea. I've chosen $x=2$. 
\begin{centred}
\verb`\solve{$ \tan x - \lambda x $}[\lambda=-1/0.8,{x}=2][5]` $\Longrightarrow$
\solve{$ \tan x - \lambda x $}[\lambda=-1/0.8,{x}=2][5].
\end{centred}
Table 4.19 of \emph{HMF }lists values of $x$ against $\lambda$ and
this is the value tabulated there.

\section{Extrema}

A function may not have a zero; or, for the given initial trial value
and initial step in the search for a zero, there may be a local extremum
in the way. In that case \verb`numerica-plus` may well locate the
local extremum (maximum or minimum but not a saddle point). For example
for the quadratic $(2x-1)^{2}+3x+1$ the \verb`\solve` command gives
the result 
\begin{centred}
\verb`\solve[vvi=]{$ (2x-1)^2+3x+1 $}[x=2]` $\Longrightarrow$\solve[vvi=]{$ (2x-1)^2+3x+1 $}[{x}=2]. 
\end{centred}
Since $(2x-1)^{2}+3x+1\ne0$ for any (real number) $x$, we deduce
that the quadratic takes a minimum value $1.9375$ at $x=0.125$ \textendash{}
easily confirmed analytically. This particular minimum is a global
minimum but in general any extremum found is only \emph{local}. The
function may well take larger or smaller values (or vanish for that
matter) further afield.

It is also worth noting in this example the \verb`vvi=` in the settings
option which suppresses display of the vv-list. (The only member of
the vv-list is the trial value \verb`x=2` which we do not want to
display.)

\noindent\begin{minipage}[t]{1\columnwidth}%
\begin{shaded}%
Note that the function for which a zero is being sought is \emph{not}
equated to zero when entered in the \verb`\solve` command. It is
\verb`\solve{ f(x) }`, not \verb`\solve{ f(x)=0 }`. This is precisely
because it may be an extremum that is found rather than a zero (if
extremum or zero is found at all \textendash{} think $e^{x}$). The
display of the result makes clear which is which, equating $f(x)$
to its value, zero or extremum depending on what has been found, as
you can see in the preceding examples.\end{shaded}%
\end{minipage}

\subsection{The search strategy}

\label{subsec:solveSearch-strategy}If you have some sense of where
a function has a zero, then choose a trial value in that vicinity.
\verb`\solve` uses a bisection method to home in on the zero. It
therefore needs \emph{two} initial values. For the first it uses the
trial value you specify, call it $a$ and for the second, by default,
it uses $a+1$. (The default value $1$ for the initial step from
the trial value can be changed in the settings option; see �\ref{sec:solveSettings-option}.)
If $f(a)$ and $f(a+1)$ have opposite signs then that is good. Bisection
of the interval $[a,a+1]$ can begin immediately in order to home
in on the precise point where $f$ vanishes. Write $b=a+1$.
\begin{itemize}
\item Let $c=\tfrac{1}{2}(a+b)$; if $f(c)=0$ the zero is found; otherwise
either $f(a),f(c)$ are of opposite signs or $f(c),f(b)$ are of opposite
signs. In the former case write $a_{1}=a,$ $b_{1}=c$; in the latter
case write $a_{1}=c$, $b_{1}=b$ and then redefine $c=\tfrac{1}{2}(a_{1}+b_{1})$.
Continue the bisection process, either until an exact zero $c$ of
$f$ is reached ($f(c)=0$) or a value $c$ is reached where the difference
between $a_{n+1}$ and $b_{n+1}$ is zero at the specified rounding
value. (But note, $f(c)$ may not vanish at that rounding value \textendash{}
the zero might be elsewhere in the interval and $f$ might cross the
axis at a steep slope.)
\end{itemize}
However $f(a)$ and $f(b)=f(a+1)$ may not have opposite signs. If
we graph the function $y=f(x)$ and suppose $f(a),f(b)$ are distinct
but of the same sign, then the line through the points $(a,f(a))$,
$(b,f(b))$ will intersect the $x$-axis to the left of $a$ or the
right of $b$ depending on its slope. We search always \emph{towards
the $x$-axis} in steps of $b-a$ ($=1$ with default values).
\begin{itemize}
\item If the line intersects the axis to the left of $a$ then $c=a-(b-a)$
and we set $a_{1}=c,b_{1}=a$; if the line intersects the axis to
the right of $b$ then $c=b+(b-a)$ and we set $b_{1}=c,a_{1}=b$.
The hope is that by always taking steps in the direction towards the
$x$-axis that eventually $f(c)$ will be found to lie on the \emph{opposite}
side of the axis from $f(a_{n})$ or $f(b_{n})$, at which point the
bisection process begins.
\item Of course this may not happen. At some point $c$ may lie to the left
of $a_{n}$ but $\left|f(c)\right|>\left|f(a_{n})\right|$, or $c$
may lie to the right of $b_{n}$ but $\left|f(c)\right|>\left|f(b_{n})\right|$.
The slope has reversed. In that case we halve the step value to $\tfrac{1}{2}(b-a)$
and try again in the same direction as before from the same point
as before ($a_{n}$ or $b_{n}$ as the case may be).
\item Should we find at some point that $f(a_{n})=f(b_{n})$ then the previous
strategy does not apply. In this case we choose $a_{n+1}$ and \textbf{$b_{n+1}$}
at the quarter and three-quarter marks between $a_{n}$ and $b_{n}$.
Either $f(a_{n+1})$ and $f(b_{n+1})$ will differ and the previous
search strategy can start again or we are on the way to finding an
extremum of $f$.
\end{itemize}
As already noted it is also possible that our function has neither
zeros nor extrema. To prevent the search continuing indefinitely,
\texttt{numerica} uses a cut-off value for the maximum number of steps
pursued \textendash{} by default set at 100.

\subsubsection{Elusive extrema}

The strategy `search always towards the $x$-axis' has a consequence:
it means that a local maximum above the $x$-axis will almost certainly
not be found, since `towards the $x$-axis' pulls the search away
from the maximum. Similarly a local minimum below the $x$-axis will
also not be found since `towards the $x$-axis' pulls the search
away from the minimum. 

One way of countering this elusiveness is to add a constant value
(possibly negative) to the function whose zeros and extrema are being
sought. The zeros of the function will change but the abscissae ($x$
values) of the extrema remain unchanged. If the constant is big enough
it will push a local minimum above the axis where it can be found
or, for a negative constant, push a local maximum below the axis where
it can be found.

For example $f(x)=x^{3}-x$ has roots at $-1,0,1$, a local maximum
at $-\tfrac{1}{\surd3}$ and a local minimum at $\tfrac{1}{\surd3}$.
To locate the minimum, I have added an unnecessarily large constant
$k$ to $f(x)$. ($k=1$ would have sufficed, but note, $k=0$ fails.)
\begin{centred}
\verb`\solve{$ x^3-x+k $}[k=5,{x}=0.5]` $\Longrightarrow$ \solve{$ x^3-x+k $}[k=5,{x}=0.5].
\end{centred}
Checking, \verb`\eval{$\tfrac1{\surd 3}$}` $\Longrightarrow$ \eval{$\tfrac1{\surd 3}$}.
There is a discrepancy in the $6$th decimal place which can be eliminated
by using the extra rounding setting; see �\ref{subsec:solveExtraRounding}.

\subsubsection{False extrema}

A function which `has an infinity' at a particular value can result
in a false extremum being found:
\begin{centred}
\verb`\solve{$ 1/x $}[x=-1/3]` $\Longrightarrow$ \solve{$ 1/x $}[x=-1/3].
\end{centred}
One needs to look for extrema with some awareness, a general sense
of how the function behaves. `Searching blind' may lead to nonsense
results. In this particular example, changing the rounding value will
show the supposed extremum jumping from one large value to another
and not settling at a particular value.

\section{Star (\texttt{{*}}) option}

A starred form of the\textbf{ }\verb`\nmcSolve` command suppresses
all elements of display of the result apart from the numerical value.
In v.1 of \verb`numerica`, when the commands of \verb`numerica-plus`
were invoked with a package option, nesting a \verb`\solve*` command
within another command was the form to use. Now that \verb`numerica-plus`
is a separate package (but labelled v.2), the star is no longer necessary.
\verb`numerica` and associated packages understand a nested command
to be the starred form, whether the star is explicitly present or
not. 

With the `elusive' extremum example above, we can find the actual
value of the minimum by nesting \verb`\solve*` or \verb`\solve`
within the vv-list of an \verb`\eval` command:
\begin{centred}
\verb`\eval{$ x^3-x $}[x={\solve{y^3-y+k}[k=5,y=0.5]}]` $\Longrightarrow$
\eval{$ x^3-x $}[x={\solve{y^3-y+k}[k=5,y=0.5]}].
\end{centred}
(Note the braces around the \verb`\solve` and its vv-list to hide
\emph{its} square-brackets  from the parsing of the vv-list of the
\verb`\eval` command.) The result is to be compared with \verb`$\eval*{x^3-x}[x=\tfrac1{\surd3}]$`
$\Longrightarrow$ $\eval*{x^3-x}[x=\tfrac1{\surd3}]$. 

\section{Settings option}

\label{sec:solveSettings-option}The settings option is a comma-separated
list of items of the form \emph{key~=~value}.

\subsection{Inherited settings}

The keys discussed in the settings\emph{ }option for \verb`\nmcEvaluate`
are also available for \verb`\nmcSolve`. The very first example in
this chapter used the punctuation option \texttt{p} (\verb`\solve[p]{\[... `)
inherited from the \verb`\eval` command to ensure a comma after the
display-style presentation of the result. We also saw in the quadratic
example illustrating extrema the use of \texttt{vvi} with no value
to suppress display of the vv-list: \verb`\solve[vvi=]{$ ...`. 

Putting \verb`dbg=1` produces a familiar kind of display. Using the
function 
\[
ct-\sqrt{a^{2}+b^{2}-2ab\cos(\beta+\omega t)}
\]
from the rotating disk problem,
\begin{verbatim}
    \solve[dbg=1,var=t]
      {$ ct-\sqrt{a^{2}+b^{2}-2ab\cos(\beta+\omega t)} 
      $}[a=10,b=20,\beta=1,{t}=0][4]
\end{verbatim}
$\Longrightarrow$     \solve[dbg=1,var=t]
      {$ ct-\sqrt{a^{2}+b^{2}-2ab\cos(\beta+\omega t)} 
      $}[a=10,b=20,\beta=1,{t}=0][4]

\subsubsection{Multi-line display of the result}

\label{subsec:solveMulti-line-display}By default the result is presented
on a single line. Unless the star option is being used, this can be
of the form \emph{function = function value, (vv-list) $\rightarrow$
result}. It takes only a slightly complicated formula and only a few
variables in the vv-list before this becomes a crowded line, likely
to exceed the line width and extend into the margin. To split the
display over two lines choose a \texttt{vvd} specification in the
vv-list like, for instance, \texttt{vvd=\{,\}\textbackslash\textbackslash (vv)}.
The \texttt{\textbackslash\textbackslash} is a trigger for \texttt{numerica}
to replace whatever environment the \texttt{\textbackslash eval}
command is wrapped around with a \texttt{multline} environment. An
asterisk in the vv-list replaces \texttt{multline} with \texttt{multline{*}}
so that no equation number is used:
\begin{verbatim}
    \solve[p=.,vvd={,}\\(vv),*]
      {$ ct-\sqrt{a^{2}+b^{2}-2ab\cos(\beta+\omega t)} 
      $}[a=10,b=20,\beta=1,{t}=0][4]
\end{verbatim}
$\Longrightarrow$ \solve[p=.,vvd={,}\\(vv),*]
      {$ ct-\sqrt{a^{2}+b^{2}-2ab\cos(\beta+\omega t)} 
      $}[a=10,b=20,\beta=1,{t}=0][4]
\begin{verbatim}

\end{verbatim}
You could introduce a third line if you wished to display the result
on a line of its own by using the spec. \texttt{vvd=\{,\}\textbackslash\textbackslash (vv)\textbackslash\textbackslash}.

The function evaluates to $-0.0007$. Is this a zero that has been
found or an extremum? To find out, the calculation needs to be carried
out to a higher rounding value which is the reason why \verb`\nmcSolve`
has an extra rounding setting; see �\ref{subsec:solveExtraRounding}
below.

\subsection{\texttt{\textbackslash nmcSolve}-specific settings}

In addition there are some settings specific to \verb`\nmcSolve`.
These are listed in Table~\ref{tab:solveSettings}.

\begin{table}
\centering{}\caption{\protect\label{tab:solveSettings}Settings for \texttt{\textbackslash nmcSolve}}
\begin{center}
\begin{tabular}{llll}
\toprule 
{\small key} & {\small type} & {\small meaning} & {\small default}\tabularnewline
\midrule
{\small\texttt{var}} & {\small token(s)} & {\small equation variable} & \tabularnewline
{\small\texttt{dvar}} & {\small real $\ne0$ } & {\small initial step size} & {\small\texttt{1}}\tabularnewline
{\small\texttt{+}} & {\small int} & {\small extra rounding} & {\small\texttt{0}}\tabularnewline
{\small\texttt{max}} & {\small int > 0} & {\small max. number of steps before cut off} & {\small\texttt{100}}\tabularnewline
{\small\texttt{reuse}} & {\small int ($\mathtt{0}/\mathtt{1}$)} & {\small form of result saved with }{\small{\small\verb`\reuse`}} & {\small\texttt{0}}\tabularnewline
\bottomrule
\end{tabular}
\par\end{center}
\end{table}


\subsubsection{Equation variable}

By default the equation variable is the \emph{rightmost} variable
in the vv-list. This may not always be convenient. A different equation
variable can be specified by entering
\begin{lyxcode}
var~=~<variable~name>
\end{lyxcode}
in the vv-list. \texttt{<variable name>} will generally be a single
character or token \textendash{} $x$, $t$, $\alpha$, $\omega$
\textendash{} but is not necessarily of this kind. Multi-token names
are perfectly acceptable (with the default multi-token setting; see
the associated document \verb`numerica.pdf` about this).

\subsubsection{Initial step size}

The vv-list contains the equation variable set to a trial value. But
\verb`\solve` needs \emph{two} initial values to begin its search
for a zero or extremum; see �\ref{subsec:solveSearch-strategy}. Ideally,
these values will straddle a zero of the function being investigated.
By default, the second trial value is $1$ more than the first: if
the equation variable is set to a trial value $a$ then the second
value defaults to $a+1$. The `$+1$' here can be changed by entering
in the settings option
\begin{lyxcode}
dvar~=~<non-zero~real~number>
\end{lyxcode}
For instance, \texttt{dvar=-1}, or \texttt{dvar=\textbackslash pi}
are two valid specifications of initial step size. The notation is
prompted by the use of expressions like $x$ and $x+dx$ for two nearby
points in calculus.

An example where the default step value is too big and a smaller one
needs to be specified is provided by Planck's radiation function (\emph{HMF
}Table 27.2),
\[
f(x)=\frac{1}{x^{5}(e^{1/x}-1)}.
\]
From the (somewhat coarse-grained) table in \emph{HMF }it is clear
that there is a maximum of about 21.2 when $x$ is a little more than
$0.2$. This is a maximum above the $x$-axis and hence `elusive'.
To find it, substract $100$ (say) from the formula and again use
the ability to nest commands to display the result. In the example,
I find in the vv-list of the \verb`\eval` command the value of $x$
which maximizes the Planck radiation function, then calculate the
maximum in the main argument of the \verb`\eval` command. Note the
\verb`dvar=0.1` in the settings option of the \verb`\solve*` command: 
\begin{verbatim}
    \eval[p=.]{\[ \frac1{x^5(e^{1/x}-1)} \]}
      [ x={ \solve*[dvar=0.1]
        { \frac1{y^5(e^{1/y}-1)}-100 }[y=0.1]
          } ]
\end{verbatim}
$\Longrightarrow$     \eval[p=.]{\[ \frac1{x^5(e^{1/x}-1)} \]}
      [ x={ \solve*[dvar=0.1]
              { \frac1{y^5(e^{1/y}-1)}-100 }[y=0.1]
          } ]

\noindent The maximum is indeed a little over $21.2$ and the $x$
value a little more than $0.2$. 

The default \verb`dvar=1` is too big for this problem. From the table
in \emph{HMF},\emph{ }$f(0.1)=4.540$ and $f(1.1)=0.419$. Thus for
$f(x)-100$ the `towards the $x$-axis' search strategy would lead
to negative values of $x$ with the default \verb`dvar` setting. 

\subsubsection{Extra rounding}

\label{subsec:solveExtraRounding}\verb`\solve` determines that a
zero or an extremum has been reached when the difference between two
successive bisection values vanishes at the specified rounding value
(the value in the final trailing optional argument of the \verb`\solve`
command; $6$ by default). If our function is $f(x)$ then $\abs{x_{n+1}-x_{n}}=0$
to the specified rounding value and $f(x_{n})$, $f(x_{n+1})$ have
opposite signs or at least one vanishes. Then (assuming $x_{n+1}>x_{n}$
and continuity) there must be a critical value $x_{c}\in[x_{n},x_{n+1}]$
such that $f(x_{c})=0$ exactly. But in general the critical value
$x_{c}$ will not coincide with $x_{n}$ or $x_{n+1}$. If $f(x)$
crosses the $x$-axis at a steep angle it may well be that although
$f(x_{c})$ vanishes to all $16$ figures, $f(x_{n})$ and $f(x_{n+1})$
do not, not even at the (generally smaller) specified rounding value.
For instance, suppose $f(x)=1000x-3000$ and that our trial value
is $x=e$:
\begin{centred}
\verb`\solve[vvi=]{$ 1000x-3000 $}[x=e][4*]` $\Longrightarrow$ \solve[vvi=]{$ 1000x-3000 $}[x=e][4*].
\end{centred}
Although the difference between successive $x$ values vanishes to
$4$ places of decimals, $f(x)$ does not, not even to $2$ places.
If we want the function to vanish at the specified rounding value
\textendash{} $4$ in the example \textendash{} then we will need
to locate the zero more precisely than that.

This is the purpose of the extra rounding key in the settings option.
Enter
\begin{lyxcode}
+~=~<integer>
\end{lyxcode}
in the settings option of the \verb`\solve` command to add \texttt{<integer>}
to the rounding value determining the conclusion of the calculation.
By default, \texttt{+=}0. 

With this option available it is easy to check that \verb`+=3` suffices
in the example to ensure that both $x$ and $f(x)$ vanish to $4$
places of decimals,
\begin{centred}
\verb`\solve[+=3]{$ 1000x-3000 $}[x=e][4*]` $\Longrightarrow$ \solve[+=3]{$ 1000x-3000 $}[x=e][4*],
\end{centred}
and that \texttt{+=2} does not, i.e., we need to locate the zero to
$4+3=7$ figures to ensure the function vanishes to $4$ figures. 

There is no need for the \texttt{<integer>} to be positive. In fact
negative values can illuminate what is going on. In the first of the
following, the display is to $10$ places but (\verb`+=-4`) the calculation
is only to $10-4=6$ places. In the second, the display is again to
$10$ places, but (\verb`+=-3`) the calculation is to $10-3=7$ places.
\begin{centred}
\verb`\solve[+=-4]{$ 1000x-3000 $}[x=e][10*]` $\Longrightarrow$
\solve[+=-4]{$ 1000x-3000 $}[x=e][10*],

\verb`\solve[+=-3]{$ 1000x-3000 $}[x=e][10*]` $\Longrightarrow$
\solve[+=-3]{$ 1000x-3000 $}[x=e][10*].
\end{centred}
Only in the second does $f(x)=1000x-3000$ vanish when rounded to
$4$ figures. 

Returning to an earlier example (�\ref{subsec:solveMulti-line-display})
in which it was not entirely clear whether a zero or an extremum had
been found, we can now resolve the confusion. Use the extra rounding
setting (and pad with zeros to emphasize the $4$-figure display by
adding an asterisk in the trailing optional argument):
\begin{verbatim}
    \solve[+=2,vvd={,}\\(vv),*]
      {$ ct-\sqrt{a^{2}+b^{2}-2ab\cos(\beta+\omega t)} 
        $}[a=10,b=20,\beta=1,{t}=0][4*]
\end{verbatim}
$\Longrightarrow$     \solve[+=2,vvd={,}\\(vv),*]
      {$ct-\sqrt{a^{2}+b^{2}-2ab\cos(\beta+\omega t)} 
        $}[a=10,b=20,\beta=1,{t}=0][4*]

\subsubsection{Maximum number of steps before cut-off}

Once two function values have been found of different sign, bisection
is guaranteed to arrive at a result. The problem is the \emph{search}
for two such values. This may not terminate \textendash{} think of
a function like $e^{x}$ which lacks both zeros and extrema. To prevent
an infinite loop, \verb`\solve` cuts off the search after $100$
steps. This cut-off value can be changed for a calculation by entering 
\begin{lyxcode}
max~=~<positive~integer>
\end{lyxcode}
in the settings option.

To illustrate, we know that $1/x$ has neither zero nor extremum,
but we do not get an infinite loop; we get an error message if we
attempt to `solve' $1/x$:
\begin{centred}
\verb`\solve{ 1/x }[x=1]` $\Longrightarrow$ \solve{ 1/x }[x=1]
\end{centred}

\subsubsection{Form of result saved by \texttt{\textbackslash nmcReuse}}

As wth \verb`\eval` and \verb`\iter` it is possible to specify to
some extent what is saved to file when using \verb`\reuse` after
a \verb`\solve` command. The form of entry in the settings option
is
\begin{lyxcode}
reuse~=~<integer>
\end{lyxcode}
If the star option is used with the \verb`\solve` command the numerical
result is the only thing saved, but in the absence of the star option,
\begin{itemize}
\item \texttt{reuse=0} saves\emph{ the form that is displayed}. For example,
if the display is of the form \emph{function = function value, (vv-list)
$\rightarrow$ result} then that is what is saved; this is the default
behaviour;
\item \texttt{reuse=1} (or any non-zero integer) saves only the numerical
result.
\end{itemize}

\subsection{Changing default values}

\begin{wraptable}[12]{o}{0.5\columnwidth}%
\centering{}\caption{\protect\label{tab:solveSettingsDefaults}Defaults for \texttt{\textbackslash nmcSolve}}
\begin{center}
\begin{tabular}{ll}
\toprule 
{\small key} & {\small default}\tabularnewline
\midrule
{\small solve-first-step} & 1\tabularnewline
{\small solve-extra-rounding} & {\small\texttt{0}}\tabularnewline
{\small solve-max-steps} & {\small\texttt{100}}\tabularnewline
{\small solve-reuse} & {\small\texttt{0}}\tabularnewline
\bottomrule
\end{tabular}
\par\end{center}\end{wraptable}%
If you wish to change the default values of the various settings for
\verb`\nmcSolve` this can be done by entering new values in a configuration
file \texttt{numerica-plus.cfg} as described in the chapter on settings
in the associated document \texttt{numeric.pdf}. The relevant keys
are listed in Table~\ref{tab:solveSettingsDefaults}, corresponding
to the \verb`dvar`, \verb`+`, \verb`max` and \verb`reuse` settings
of the \verb`\solve` command. (Obviously it makes no sense to have
a default setting for the solution variable. That will change from
case to case.)

\chapter{Recurrence relations: \texttt{\textbackslash nmcRecur}}

One of the simplest recurrence relations is that determining the Fibonacci
numbers, $f_{n+2}=f_{n+1}+f_{n}$, with initial values $f_{0}=f_{1}=1$.
The command \verb`\nmcRecur`, short-name form \verb`\recur`, allows
calculation of the terms of this sequence:
\begin{verbatim}
    $ \nmcRecur[do=8,see1=8,...]
      { f_{n+2}=f_{n+1}+f_{n} }
        [f_{1}=1,f_{0}=1] $
\end{verbatim}
$\Longrightarrow$     $\nmcRecur[do=8,see1=8,...]
      { f_{n+2}=f_{n+1}+f_{n} }
        [f_{1}=1,f_{0}=1]$

The recurrence relation is entered in the main argument (between braces),
the initial values in the vv-list trailing the main argument, and
the display specification is placed in the settings option: \texttt{do=8}
terms to be calculated, all $8$ to be viewed (\texttt{see1=8}), and
the display to be concluded by an ellipsis to indicate that the sequence
continues (those are three dots/periods/full stops in the settings
option, not an ellipsis glyph).

A more complicated recurrence relation determines the Legendre polynomials:
\[
(n+2)P_{n+2}(x)-(2n+3)xP_{n+1}(x)+(n+1)P_{n}(x)=0.
\]
For the purposes of \verb`\recur` we need $P_{n+2}$ expressed in
terms of the lower order terms:
\[
P_{n+2}(x)=\frac{1}{n+2}\left((2n+3)xP_{n+1}(x)-(n+1)P_{n}(x)\right).
\]
 It is this standard form \textendash{} the term to be calculated
on the left, equated to an expression involving a fixed number of
lower-order terms on the right \textendash{} that \verb`numerica-plus`
works with. For $P_{0}(x)=1,~P_{1}(x)=x$ and $x=0.5$, the terms
are calculated thus:
\begin{verbatim}
    \recur[p,do=11,see1=4,see2=2,vvd={,}\\(vv)\\,*]
      {\[ P_{n+2}(x)=\frac{1}{n+2}
        \Bigl((2n+3)xP_{n+1}(x)-(n+1)P_{n}(x)\Bigr)
      \]}[P_{1}(x)=x,P_{0}(x)=1,x=0.5]
\end{verbatim}
$\Longrightarrow$     \recur[p,do=11,see1=4,see2=2,vvd={,}\\(vv)\\,*]
      {\[ P_{n+2}(x)=\frac{1}{n+2}
        \Bigl((2n+3)xP_{n+1}(x)-(n+1)P_{n}(x)\Bigr)
      \]}[P_{1}(x)=x,P_{0}(x)=1,x=0.5] 

\noindent where $P_{9}(0.5)$ and $P_{10}(0.5)$ are the last two
displayed values (and to $6$-figures are the values listed in \emph{HMF
}Table 8.1).

These examples also illustrate a common behaviour of the commands
in \verb`numerica` and associated packages: when wrapped around math
delimiters: the display is of the \emph{expression=result} form, and
when placed between math delimiters the display is simply of the \emph{result}.
When used without math delimiters, \verb`numerica-plus` treats the
command as if it had been placed between \texttt{\textbackslash{[}
\textbackslash{]}}.

\section{Notational niceties}

More than the other commands in {\ttfamily\verb`numerica`}
and associated packages, \verb`\nmcRecur` depends on getting the
notation into a standard form.
\begin{itemize}
\item The terms of the recurrence must be \emph{subscripted}: $f_{n}$,
$P_{n}(x)$ are examples.
\item The recurrence relation is placed in the main (mandatory) argument
of \verb`\nmcRecur` in the form: \emph{high-order term=function of
lower-order terms}.
\item The initial-value terms in the vv-list must occur left-to-right in
the order \emph{high }to \emph{low} order.
\item The recurrence variable changes by $1$ between successive terms.
\end{itemize}
The example for Legendre polynomials in particular shows what is required.
The Fibonacci example is simpler, since the recurrence variable does
not occur independently in the recurrence relation as it does with
the Legendre polynomials. In both cases though the recurrence variable
is absent from the vv-list.

\subsection{Vv-list and recurrence variable}

The recurrence variable is required in the vv-list only when an implicit
mode calculation is undertaken. Suppose we write $A$ and $B$ for
the coefficients $2n+3$ and $n+1$ respectively in the Legendre recurrence.
$A$ and $B$ will now need entries in the vv-list which means the
recurrence variable will need a value assigned to it there too, and
we will need to add \verb`vv@=1` (or \verb`vvmode=1`) to the settings
option.
\begin{verbatim}
    \recur[p,vvmode=1,do=11,see1=4,see2=2,vvd={,}\\(vv)\\,*]
      {\[ P_{n+2}(x)=\frac{1}{n+2}
        \Bigl(AxP_{n+1}(x)-BP_{n}(x)\Bigr)
      \]}[P_{1}(x)=x,P_{0}(x)=1,x=0.5,A=2n+3,B=n+1,n=0]
\end{verbatim}
$\Longrightarrow$     \recur[p,vvmode=1,do=11,see1=4,see2=2,vvd={,}\\(vv)\\,*]
      {\[ P_{n+2}(x)=\frac{1}{n+2}
        \Bigl(AxP_{n+1}(x)-BP_{n}(x)\Bigr)
      \]}[P_{1}(x)=x,P_{0}(x)=1,x=0.5,A=2n+3,B=n+1,n=0] 

Since the vv-list is evaluated from the right, the left-to-right high-to-low
ordering of the initial-value terms means the value of the lowest
order term is read first. Although \verb`numerica-plus` depends on
this order of occurrence of the terms, they do not need to be \emph{consecutive}
as in the examples so far (although it is natural to enter them in
this way). \verb`numerica-plus` reads the value of the subscript
of only the right-most term (the lowest order term), increments it
by $1$ when reading the next recurrence term to the left, and so
on. The reading of the subscript of the lowest order term in the vv-list
provides the initial value of the recurrence variable.

In the following example I have placed other items between $P_{1}(x)$
and $P_{0}(x)$ in the vv-list (but maintained their left-to-right
order) and given the recurrence variable $n$ a ridiculous initial
value $\pi^{2}/12$. (Because of the order in which things get done
`behind the scenes', \emph{some} value is necessary so that the
$n$ in `$B=n+1$' does not generate an `unknown token' message.)
The result is unchanged.
\begin{verbatim}
    \recur[p,vvmode=1,do=11,see1=4,see2=2,vvd={,}\\(vv)\\,*]
      {\[ P_{n+2}(x)=\frac{1}{n+2}
        \Bigl(AxP_{n+1}(x)-BP_{n}(x)\Bigr)
      \]}[A=2n+3,P_{1}(x)=x,B=n+1,n=\pi^2/12,P_{0}(x)=1,x=0.5]
\end{verbatim}
$\Longrightarrow$     \recur[p,vvmode=1,do=11,see1=4,see2=2,vvd={,}\\(vv)\\,*]
      {\[ P_{n+2}(x)=\frac{1}{n+2}
        \Bigl(AxP_{n+1}(x)-BP_{n}(x)\Bigr)
      \]}[A=2n+3,P_{1}(x)=x,B=n+1,n=\pi^2/12,P_{0}(x)=1,x=0.5]

\subsection{Form of the recurrence relation}

As noted earler, the form of the recurrence must be entered in the
main argument in the form: \emph{highest order term = function of
consecutive lower order terms}. The number of lower\emph{ }order terms
is the order of the recurrence. The Fibonacci and Legendre polynomial
recurrences are both second order and presented in the form: \emph{$n+2$-th
term = function of $n+1$-th term and $n$-th term}. We could equally
have done
\begin{verbatim}
    \nmcRecur[p,do=8,see1=8,...]
      {$ f_{n}=f_{n-1}+f_{n-2} $}
        [f_{1}=1,f_{0}=1]
\end{verbatim}
$\Longrightarrow$     \nmcRecur[p,do=8,see1=8,...]
      {$ f_{n}=f_{n-1}+f_{n-2} $}
        [f_{1}=1,f_{0}=1] where now the recurrence is of the form $n$\emph{-th term = function
of $n-1$-th term and $n-2$-th term}, or (adjusting the coefficients
as well as the recurrence terms),
\begin{verbatim}
    \recur[p=.,do=10,see1=4,see2=2,vvd={,}\\(vv)\\,*]{\[
      P_{n+1}(x)=\frac{1}{n+1}
        \Bigl((2n+1)xP_{n}(x)-nP_{n-1}(x)\Bigr)
          \]}[P_{2}(x)=-0.125,P_{1}(x)=x,x=0.5]
\end{verbatim}
$\Longrightarrow$     \recur[p=.,do=10,see1=4,see2=2,vvd={,}\\(vv)\\,*]{\[
      P_{n+1}(x)=\frac{1}{n+1}
        \Bigl((2n+1)xP_{n}(x)-nP_{n-1}(x)\Bigr)
          \]}[P_{2}(x)=-0.125,P_{1}(x)=x,x=0.5]

\noindent The recurrence here is of the form $n+1$\emph{-th term
= function of $n$-th term and $n-1$-th term}. This last example
has one further `wrinkle'. I've made $P_{1}(x)$ the lowest order
term and decreased the number of terms to calculate by $1$ accordingly. 

\subsection{First order recurrences (iteration)}

The recurrence relations for both the Fibonacci sequence and Legendre
polynomials are second order. There is no reason why the recurrence
should not be of third or higher order or, indeed, lower. A first
order recurrence provides an alternative means of iterating functions.
\verb`\recur` therefore provides a means to display the results of
an iteration in a different form from \verb`\iter`. 

Iterating $1+a/x$ in this way, $16$ terms gives the sequence
\begin{verbatim}
    \recur[do=16,see1=0,see2=3,...]{$
      x_{n+1}=1+a/x_{n}
        $}[x_{0}=1,a=1]
\end{verbatim}
$\Longrightarrow$     \recur[do=16,see1=0,see2=3,...]{$
      x_{n+1}=1+a/x_{n} 
        $}[x_{0}=1,a=1]

\noindent to be compared with the example near the start of Chapter~\ref{chap:Iterating-functions}.
(\emph{That} effected $15$ iterations; \emph{this} uses $16$ terms
because of the extra $x_{0}=1$ term.) 

\section{Star (\texttt{{*}}) option}

When the star option is used with the \verb`\nmcRecur` command, only
a single term, the \emph{last}, is presented as the result. Repeating
the last calculation, but with the star option produces
\begin{verbatim}
    \recur*[p=.,do=10]{\[
      P_{n+1}(x)=\frac{1}{n+1}
        \Bigl((2n+1)xP_{n}(x)-nP_{n-1}(x)\Bigr)
          \]}[P_{2}(x)=-0.125,P_{1}(x)=x,x=0.5]
\end{verbatim}
$\Longrightarrow$     \recur*[p=.,do=10]{\[
      P_{n+1}(x)=\frac{1}{n+1}
        \Bigl((2n+1)xP_{n}(x)-nP_{n-1}(x)\Bigr)
          \]}[P_{2}(x)=-0.125,P_{1}(x)=x,x=0.5]

Although punctuation (a full stop) was specified in the settings,
it has been ignored in the display of the result. Other settings would
also have been ignored with the exception of the \verb`do` key which
is required to know exactly which term to calculate. The star option
produces a purely numerical answer without any trimmings.

\noindent\begin{minipage}[t]{1\columnwidth}%
\begin{shaded}%
This seems something of a waste of the star option since it gives
much the same result as choosing \texttt{do=10,see1=0,see2=1}. Not
\emph{exactly} the same, since math delimiters are involved now, but
sufficiently similar to make me wonder if I should change the starred
form to apply only to those recurrences which approach a limit. The
starred form would then produce the limiting value as its result (like
\verb`\iter*`). This is a possible change for future versions of
\verb`numerica-plus` and should be borne in mind if using \verb`\recur*`.\end{shaded}%
\end{minipage}

\section{Settings}

The settings option is a comma-separated list of items of the form
\emph{key~=~value}.

\subsection{Inherited settings}

Because recurrence terms are necessarily multi-token, the multi-token
key is hard-coded in \verb`\recur` to \texttt{xx=1}. 

\subsubsection{Multi-line formatting of result}

When the \verb`\recur` command wraps around math delimiters, the
\texttt{vvd} setting is available to split display of the result over
two or more lines. For example, \texttt{vvd=\{,\}\textbackslash\textbackslash (vv)}pushes
the vv-list and sequence of calculated values to a second line; or,
\texttt{vvd=\{,\}\textbackslash qquad(vv)\textbackslash\textbackslash}
pushes only the sequence of calculated values to a second line; or
\texttt{vvd=\{,\}\textbackslash\textbackslash (vv)\textbackslash\textbackslash}
pushes the vv-list, centred, to a second line and the sequence of
values, right aligned, to a third line. The \texttt{{*}} setting is
available to suppress equation numbering (by substituting \verb`multline*`
for \verb`multline`).
\begin{verbatim}
 \nmcRecur[do=8,see1=8,...,vvd={,}\qquad(vv)\\,*]
      {$ f_{n+2}=f_{n+1}+f_{n} $}
        [f_{1}=1,f_{0}=1]
\end{verbatim}
$\Longrightarrow$     \nmcRecur[do=8,see1=8,...,vvd={,}\qquad(vv)\\,*]
      {$ f_{n+2}=f_{n+1}+f_{n} $}
        [f_{1}=1,f_{0}=1]

\subsection{\texttt{\textbackslash nmcRecur}-specific settings}

\label{subsec:recurSpecific-settings}In addition to the inherited
settings there are some specific to \verb`\nmcRecur`. These are listed
in Table~\ref{tab:recurSettings} below.

\subsubsection{Number of terms to calculate}

By entering
\begin{lyxcode}
do~=~<integer>
\end{lyxcode}
in the settings option you can specify how many terms of a recurrence
to calculate. The default is set to $7$ (largely to show a sufficient
number of terms of the Fibonacci series to begin to be interesting).
Note that \texttt{<integer>} will generally \emph{not} correspond
to the subscript on the last term calculated since that also depends
on the value of the subscript of the lowest order term in the vv-list.

\subsubsection{Number of terms to display}

\begin{table}[b]
\centering{}\caption{\protect\label{tab:recurSettings}Settings for \texttt{\textbackslash nmcRecur}}
\begin{center}
\begin{tabular}{llll}
\toprule 
{\small key} & {\small type} & {\small meaning} & {\small default}\tabularnewline
\midrule
{\small\texttt{do}} & {\small int$\ge0$} & {\small number of terms to calculate} & {\small\texttt{7}}\tabularnewline
{\small\texttt{see1}} & {\small int$\ge0$} & {\small number of initial terms to display} & {\small\texttt{3}}\tabularnewline
{\small\texttt{see2}} & {\small int$\ge0$} & {\small number of final terms to display} & {\small\texttt{2}}\tabularnewline
{\small\texttt{...}} & {\small chars} & {\small follow display of values with an ellipsis} & \tabularnewline
{\small\texttt{reuse}} & {\small int ($\mathtt{0}/\mathtt{1}/\mathtt{2}$)} & {\small form of result saved with }{\small{\small\verb`\reuse`}} & {\small\texttt{0}}\tabularnewline
\bottomrule
\end{tabular}
\par\end{center}
\end{table}
By entering
\begin{lyxcode}
see1~=~<integer1>,~see2=<integer2>
\end{lyxcode}
in the settings option, you can specify how many initial terms of
the recurrence and how many of the final terms calculated you want
to view. If the sum of these settings is less than the \texttt{do}
setting, then the terms are displayed with an intervening ellipsis.
If the sum is greater than the \texttt{do} setting, then the values
are adjusted so that their sum equals the \texttt{do} setting and
all terms are displayed. 

The adjustment is preferentially to \texttt{see1}. Suppose \texttt{do=}7,
\texttt{see1=5}, \texttt{see2=4}. Then \texttt{see2} is left unchanged
but \texttt{see1} is reduced to \texttt{7-4=3}. If, say, \texttt{do=}7,
\texttt{see1=5}, \texttt{see2=8}, then \texttt{see2} is reduced to
\texttt{7} and \texttt{see1} to \texttt{-1} (rather than zero, for
technical reasons). The reason for preserving \texttt{see2} over \texttt{see1}
is for the functioning of the \texttt{reuse} setting (see below).

The default value for \texttt{see1} is $3$; the default value for
\texttt{see2} is $2$.

\subsubsection{Ellipsis}

Including three dots in the settings option
\begin{lyxcode}
...
\end{lyxcode}
ensures that a (proper) ellipsis is inserted after the final term
is displayed. An example is provided by the display of the Fibonacci
sequence at the start of this chapter. By default this option is turned
off.

\subsubsection{Form of result saved by \texttt{\textbackslash nmcReuse}}

By entering
\begin{lyxcode}
reuse~=~<integer>
\end{lyxcode}
it is possible to specify the form of result that is saved when using
\verb`\nmcReuse`. (This setting has no effect when the star option
is used with \verb`\nmcRecur`. In that case only the numerical result
of the final term calculated is saved.) There are three different
outputs possible:
\begin{itemize}
\item \texttt{int=0} (or any integer $\mathtt{\ne1,2}$) saves the full
display (the default);
\item \texttt{int=1} saves a comma-separated list of braced pairs of the
form: \texttt{\{$k$, value-of-term-$k$\}} for the last \texttt{see2}
terms calculated;
\item \texttt{int=2} saves a comma-separated list of the values of the last
\texttt{see2} terms calculated.
\end{itemize}
As an example, using \texttt{reuse=1},
\begin{verbatim}
    \recur[reuse=1,p=.,vvmode=1,do=11,see1=4,see2=2,
      vvd={,}\\(vv)\\,*]
      {\[ P_{n+2}(x)=\frac{1}{n+2} 
            \Bigl(kxP_{n+1}(x)-(n+1)P_{n}(x)\Bigr)
      \]}[k=2n+3,n=1,P_{1}(x)=x,P_{0}(x)=1,x=0.5]
    \reuse[legendre]
\end{verbatim}
$\Longrightarrow$     \recur[reuse=1,p=.,vvmode=1,do=11,see1=4,see2=2,
      vvd={,}\\(vv)\\,*]
      {\[ P_{n+2}(x)=\frac{1}{n+2} 
            \Bigl(kxP_{n+1}(x)-(n+1)P_{n}(x)\Bigr)
      \]}[k=2n+3,n=0,P_{1}(x)=x,P_{0}(x)=1,x=0.5]
    \reuse[renew]{ legendre}

\noindent Now check to see what has been saved: 
\begin{centred}
\verb`$\legendre$` $\Longrightarrow$ $ \legendre $.
\end{centred}
As you can see, the final two (because of \texttt{see2=2}) of the
$11$ Legendre polynomials calculated ($P_{0}(x)$ is the first) have
been saved, each value preceded by its index value. If the setting
had been \texttt{reuse=2}, only the two values without the index values
would have been saved. To see that they are saved as braced pairs,
use \TeX 's \verb`\meaning` command:
\begin{verbatim}
    \meaning \legendre
\end{verbatim}
$\Longrightarrow$ \meaning \legendre

\subsection{Changing default values}

\begin{wraptable}[12]{o}{0.5\columnwidth}%
\centering{}\caption{\protect\label{tab:recurSettingsDefaults}Defaults for \texttt{\textbackslash nmcRecur}}
\begin{center}
\begin{tabular}{ll}
\toprule 
{\small key} & {\small default}\tabularnewline
\midrule
{\small recur-do} & {\small\texttt{7}}\tabularnewline
{\small recur-see-first} & {\small\texttt{3}}\tabularnewline
{\small recur-see-last} & {\small\texttt{2}}\tabularnewline
{\small recur-reuse} & {\small\texttt{0}}\tabularnewline
\bottomrule
\end{tabular}
\par\end{center}\end{wraptable}%
If you wish to change the default values of the various settings for
\verb`\nmcRecur` this can be done by entering new values in a configuration
file \texttt{numerica-plus.cfg} as described in the chapter on settings
in the associated document \texttt{numerica.pdf}. The relevant keys
are listed in Table~\ref{tab:recurSettingsDefaults}, corresponding
to the \verb`do`, \verb`see1`, \verb`see2` and \verb`reuse` settings
of the \verb`\recur` command. 

\subsection{Orthogonal polynomials}

I've used Legendre polynomials in examples above, but orthogonal polynomials
generally lend themselves to the \verb`\recur` treatment. Quoting
from \emph{HMF} 22.7, orthogonal polynomials $f_{n}$ satisfy recurrence
relations of the form
\[
a_{1n}f_{n+1}(x)=(a_{2n}+a_{3n}x)f_{n}(x)-a_{4n}f_{n-1}(x),
\]
or in the standard form required by \verb`\recur`,
\[
f_{n+1}(x)=\frac{a_{2n}+a_{3n}x}{a_{1n}}f_{n}(x)-\frac{a_{4n}}{a_{1n}}f_{n-1}(x).
\]
\emph{HMF} 22.7 provides a listing of the coefficients $a_{in}$ for
the polynomials of Jacobi, Chebyshev, Legendre, Laguerre, Hermite
and others, and tables for these polynomials.

For example, Laguerre polynomials satisfy the recurrence 
\[
L_{n+1}(x)=\frac{2n+1-x}{n+1}L_{n}(x)-\frac{n}{n+1}L_{n-1}(x).
\]
with initial values $L_{0}(x)=1$ and $L_{1}(x)=1-x$. So let's calculate
the first $13$ Laguerre polynomials for, say, $x=0.5$: 
\begin{verbatim}
    \recur[do=13,see1=4,see2=2,vvd={,}\\(vv)\\,*]{\[
      L_{n+1}(x)=\frac{2n+1-x}{n+1}L_{n}(x)-
        \frac{n}{n+1}L_{n-1}(x) 
           \]}[L_{1}(x)=1-x,L_{0}(x)=1,x=0.5]
\end{verbatim}
$\Longrightarrow$     \recur[do=13,see1=4,see2=2,vvd={,}\\(vv)\\,*]{\[
      L_{n+1}(x)=\frac{2n+1-x}{n+1}L_{n}(x)-
        \frac{n}{n+1}L_{n-1}(x) 
           \]}[L_{1}(x)=1-x,L_{0}(x)=1,x=0.5]

\noindent and for $x=5$: 
\begin{verbatim}
    \recur[p=.,do=13,see1=4,see2=2,vvd={,}\\(vv)\\,*]{\[
      L_{n+1}(x)=\frac{2n+1-x}{n+1}L_{n}(x)-
        \frac{n}{n+1}L_{n-1}(x) 
           \]}[L_{1}(x)=1-x,L_{0}(x)=1,x=5]
\end{verbatim}
$\Longrightarrow$    \recur[p=.,do=13,see1=4,see2=2,vvd={,}\\(vv)\\,*]{\[
      L_{n+1}(x)=\frac{2n+1-x}{n+1}L_{n}(x)-
        \frac{n}{n+1}L_{n-1}(x) 
           \]}[L_{1}(x)=1-x,L_{0}(x)=1,x=5]

\noindent The results (reassuringly) coincide with those provided
in \emph{HMF }Table 22.11.

\subsection{Nesting}

It is possible to use the \verb`\recur` command (with star in v.1
of \verb`numerica`, with or without star in v.2 of \verb`numerica-plus`)
within an \verb`\eval`, \verb`\iter`, or \verb`\solve` command,
and indeed in \verb`\recur` itself, but with this caveat: if \verb`\recur`
is nested within another command, the initial terms of the recurrence
\textendash{} e.g., $f_{1}=1,f_{0}=1$, for the Fibonacci series,
or $L_{1}(x)=1-x,L_{0}(x)=1$ for the Laguerre polynomials \textendash{}
\emph{must be located in the vv-list of that inner }\verb`\recur`\emph{
command}. Other shared variables can often be shifted to the vv-list
of the outer command, but not these initial terms.

In the following example I multiply together (rather futilely) the
third and fourth members of the sequence of Laguerre polynomials for
$x=5$ (the answer expected is \verb`$ \eval{3.5\times2.666667} $`
$\Longrightarrow$ $ \eval{3.5\times2.666667} $). Note that although
it is tempting to shift the shared vv-lists of the inner \verb`\recur*`
commands to the vv-list of the outer \verb`\eval` command, in fact
only the \verb`x=5` entry has been transferred:
\begin{verbatim}
    \eval[p=.]{$
      \recur*[do=3]
        { L_{n+1}(x)=\frac{2n+1-x}{n+1}L_{n}(x)-
            \frac{n}{n+1}L_{n-1}(x)}
          [L_{1}(x)=1-x,L_{0}(x)=1]
      \times 
      \recur*[do=4]
        { L_{n+1}(x)=\frac{2n+1-x}{n+1}L_{n}(x)-
            \frac{n}{n+1}L_{n-1}(x)}
          [L_{1}(x)=1-x,L_{0}(x)=1]
    $}[x=5]
\end{verbatim}
$\Longrightarrow$     \eval[p=.]{$
      \recur*[do=3]
        { L_{n+1}(x)=\frac{2n+1-x}{n+1}L_{n}(x)-
            \frac{n}{n+1}L_{n-1}(x)}
              [L_{1}(x)=1-x,L_{0}(x)=1,x=5]
      \times 
      \recur*[do=4]
        { L_{n+1}(x)=\frac{2n+1-x}{n+1}L_{n}(x)-
            \frac{n}{n+1}L_{n-1}(x)}
              [L_{1}(x)=1-x,L_{0}(x)=1,x=5]
    $}

\noindent\begin{minipage}[t]{1\columnwidth}%
\begin{shaded}%
The terms of a recurrence relation are multi-token variables but \texttt{numerica}
requires single tokens for its calculations. The problem for \verb`\recur`
is that the terms in the recurrence relation in the main (mandatory)
argument differ from the terms in the vv-list: for instance $f_{n}$
in the main argument, $f_{0}$ in the vv-list. If left like that,
when \texttt{numerica} does its conversion from multi-token to single
token variables, $f_{n}$ would not be found since it differs from
$f_{0}$. Hence a crucial first step for \verb`\recur` is to reconcile
the different forms, which it does by converting the forms in the
vv-list to the forms in the recurrence in the main argument. To be
available for this form change, they must reside in the \emph{inner}
vv-list. In the outer vv-list they would be inaccessible to the inner
command.

{*}{*}{*}

This suggests an alternative way of proceeding: write the inital values
of the recurrence terms in the \emph{same} form in which they occur
in the recurrence relation, together with an initial value for the
recurrence variable, e.g., $f_{n+1}=1,f_{n}=1,n=0$. This is not how
mathematicians write the initial values in recurrence relations, which
is why I did not pursue it, but it neatly sidesteps what is otherwise
an initial awkwardness. \end{shaded}%
\end{minipage}

\chapter{Reference summary}

\section{Commands defined in \texttt{numerica-plus}}
\begin{enumerate}
\item \texttt{\textbackslash nmcIterate, \textbackslash iter}
\item \texttt{\textbackslash nmcSolve, \textbackslash solve}
\item \textbackslash\texttt{nmcRecur, \textbackslash recur}
\end{enumerate}

\section{Settings for the three commands}

\subsection{Settings for \texttt{\textbackslash nmcIterate}}

Settings option of \verb`\nmcIterate`:

\begin{center}
\begin{tabular}{llll}
\toprule 
{\small key} & {\small type} & {\small meaning} & {\small default}\tabularnewline
\midrule
{\small\texttt{var}} & {\small token(s)} & {\small iteration variable} & \tabularnewline
{\small\texttt{+}} & {\small int} & {\small fixed point extra rounding} & {\small\texttt{0}}\tabularnewline
{\small\texttt{max}} & {\small int > 0} & {\small max. iteration count (fixed points)} & {\small\texttt{100}}\tabularnewline
{\small\texttt{do}} & {\small int > 0} & {\small number of iterations to perform} & {\small\texttt{5}}\tabularnewline
{\small\texttt{see}} & {\small int > 0} & {\small number of final iterations to view} & {\small\texttt{4}}\tabularnewline
{\small\texttt{reuse}} & {\small int ($\mathtt{0}/\mathtt{1}/\mathtt{2}$)} & {\small form of result saved with }{\small{\small\verb`\reuse`}} & {\small\texttt{0}}\tabularnewline
\bottomrule
\end{tabular}
\par\end{center}

Configuration settings for \verb`\nmcIterate`:

\begin{center}
\begin{tabular}{ll}
\toprule 
{\small key} & {\small default}\tabularnewline
\midrule
{\small iter-extra-rounding} & {\small\texttt{0}}\tabularnewline
{\small iter-max-iterations} & {\small\texttt{100}}\tabularnewline
{\small iter-do} & {\small\texttt{5}}\tabularnewline
{\small iter-see-last} & {\small\texttt{4}}\tabularnewline
{\small iter-reuse} & {\small\texttt{0}}\tabularnewline
\bottomrule
\end{tabular}
\par\end{center}

\subsection{Settings for \texttt{\textbackslash nmcSolve}}

Settings option of \verb`\nmcSolve`:

\begin{center}
\begin{tabular}{llll}
\toprule 
{\small key} & {\small type} & {\small meaning} & {\small default}\tabularnewline
\midrule
{\small\texttt{var}} & {\small token(s)} & {\small equation variable} & \tabularnewline
{\small\texttt{dvar}} & {\small real $\ne0$} & {\small initial step size} & {\small\texttt{1}}\tabularnewline
{\small\texttt{+}} & {\small int} & {\small extra rounding} & {\small\texttt{0}}\tabularnewline
{\small\texttt{max}} & {\small int > 0} & {\small max. number of steps before cut off} & {\small\texttt{100}}\tabularnewline
{\small\texttt{reuse}} & {\small int ($\mathtt{0}/\mathtt{1}$)} & {\small form of result saved with }{\small{\small\verb`\reuse`}} & {\small\texttt{0}}\tabularnewline
\bottomrule
\end{tabular}
\par\end{center}

Configuration settings for \verb`\nmcSolve`:

\begin{center}
\begin{tabular}{ll}
\toprule 
{\small key} & {\small default}\tabularnewline
\midrule
{\small solve-first-step} & 1\tabularnewline
{\small solve-extra-rounding} & {\small\texttt{0}}\tabularnewline
{\small solve-max-steps} & {\small\texttt{100}}\tabularnewline
{\small solve-reuse} & {\small\texttt{0}}\tabularnewline
\bottomrule
\end{tabular}
\par\end{center}

\subsection{Settings for \texttt{\textbackslash nmcRecur}}

Settings option of \verb`\nmcRecur`:

\begin{center}
\begin{tabular}{llll}
\toprule 
{\small key} & {\small type} & {\small meaning} & {\small default}\tabularnewline
\midrule
{\small\texttt{do}} & {\small int$\ge0$} & {\small number of terms to calculate} & {\small\texttt{7}}\tabularnewline
{\small\texttt{see1}} & {\small int$\ge0$} & {\small number of initial terms to display} & {\small\texttt{3}}\tabularnewline
{\small\texttt{see2}} & {\small int$\ge0$} & {\small number of final terms to display} & {\small\texttt{2}}\tabularnewline
{\small\texttt{...}} & {\small chars} & {\small follow display of values with an ellipsis} & \tabularnewline
{\small\texttt{reuse}} & {\small int ($\mathtt{0}/\mathtt{1}/\mathtt{2}$)} & {\small form of result saved with }{\small{\small\verb`\reuse`}} & {\small\texttt{0}}\tabularnewline
\bottomrule
\end{tabular}
\par\end{center}

Configuration settings for \verb`\nmcRecur`:

\begin{center}
\begin{tabular}{ll}
\toprule 
{\small key} & {\small default}\tabularnewline
\midrule
{\small recur-do} & {\small\texttt{7}}\tabularnewline
{\small recur-see-first} & {\small\texttt{3}}\tabularnewline
{\small recur-see-last} & {\small\texttt{2}}\tabularnewline
{\small recur-reuse} & {\small\texttt{0}}\tabularnewline
\bottomrule
\end{tabular}
\par\end{center}
\end{document}
